\documentclass[11pt,a4paper]{article}

% === Typography ===
\usepackage[utf8]{inputenc}
\usepackage[T1]{fontenc}
\usepackage{mathptmx}           % Times Roman for text and math
\usepackage[scaled=0.92]{helvet} % Helvetica for sans-serif
\usepackage{courier}             % Courier for monospace
\usepackage{microtype}           % Improved kerning and line breaks

% === Layout ===
\usepackage[margin=1in]{geometry}
\usepackage{setspace}
\setstretch{1.08}                % Slightly more than single spacing
\usepackage{parskip}             % Paragraph spacing instead of indent
\setlength{\parskip}{0.5\baselineskip}
\setlength{\parindent}{0pt}

% === Mathematics ===
\usepackage{amsmath,amssymb,amsthm}

% === Tables and Figures ===
\usepackage{graphicx}
\usepackage{booktabs}
\usepackage{multirow}
\usepackage{caption}
\usepackage{subcaption}
\captionsetup{
  font={small},
  labelfont={bf},
  labelsep=period,
  skip=6pt
}

% === Algorithms ===
\usepackage{algorithm}
\usepackage{algorithmic}

% === Colors and Links ===
\usepackage{xcolor}
\definecolor{linkblue}{RGB}{0,51,153}
\definecolor{citegreen}{RGB}{0,102,51}
\usepackage[
  colorlinks=true,
  linkcolor=linkblue,
  citecolor=citegreen,
  urlcolor=linkblue,
  bookmarks=true,
  bookmarksnumbered=true,
  pdfstartview=FitH,
  pdftitle={S(M,T): A Unified Scoring Framework for LLM Routing with Provable Convergence and the Measurement Gap},
  pdfauthor={Pranav Lakherwal}
]{hyperref}

% === Bibliography ===
\usepackage{natbib}
\bibliographystyle{plainnat}

% === Headers and Footers ===
\usepackage{fancyhdr}
\pagestyle{fancy}
\fancyhf{}
\fancyhead[L]{\small\textit{S(M,T): Unified LLM Routing}}
\fancyhead[R]{\small\thepage}
\renewcommand{\headrulewidth}{0.4pt}
\fancypagestyle{plain}{%
  \fancyhf{}
  \fancyfoot[C]{\small\thepage}
  \renewcommand{\headrulewidth}{0pt}
}

% === Section Formatting ===
\usepackage{titlesec}
\titleformat{\section}{\large\bfseries}{\thesection}{1em}{}
\titleformat{\subsection}{\normalsize\bfseries}{\thesubsection}{1em}{}
\titleformat{\subsubsection}{\normalsize\itshape}{\thesubsubsection}{1em}{}
\titlespacing*{\section}{0pt}{1.5\baselineskip}{0.75\baselineskip}
\titlespacing*{\subsection}{0pt}{1\baselineskip}{0.5\baselineskip}
\titlespacing*{\subsubsection}{0pt}{0.75\baselineskip}{0.25\baselineskip}

% === Theorem Environments ===
\newtheorem{theorem}{Theorem}[section]
\newtheorem{proposition}[theorem]{Proposition}
\newtheorem{lemma}[theorem]{Lemma}
\newtheorem{corollary}[theorem]{Corollary}
\newtheorem{definition}[theorem]{Definition}
\newtheorem{remark}[theorem]{Remark}

% === Custom Commands ===
\newcommand{\smt}{S(M,T)}
\newcommand{\phii}{\phi_i}
\newcommand{\betavec}{\boldsymbol{\beta}}
\newcommand{\wvec}{\mathbf{w}}
\newcommand{\R}{\mathbb{R}}
\newcommand{\E}{\mathbb{E}}
\newcommand{\Prob}{\mathbb{P}}

% =====================================================================
\begin{document}

% === Title Block ===
\begin{center}
{\LARGE\bfseries S(M,T): A Unified Scoring Framework for LLM Routing\\[4pt]
with Provable Convergence and the Measurement Gap\par}

\vspace{14pt}

{\large Pranav Lakherwal}\\[4pt]
{\normalsize\texttt{pranav@pranavlakherwal.com}}

\vspace{10pt}

{\normalsize March 2026}
\end{center}

\vspace{10pt}

% === Abstract ===
\begin{abstract}
\noindent
We present S(M,T), a unified scoring framework for routing queries to optimal language model endpoints. The framework combines multiplicative gating, weighted compatibility functions, and Boltzmann cost penalties into a single equation: $\smt = \left[\prod_j g_j(M,T)\right] \cdot \left[\sum_i w_i \cdot \phii(M,T)\right] \cdot e^{-\betavec(T) \cdot \mathbf{H}}$. We prove convergence under Hajek conditions with logarithmic cooling, establish $O(\sqrt{KT\log T})$ regret bounds via contextual bandits reduction, and demonstrate weight identifiability when the phi matrix has full column rank. Grounding the framework against 2.76 million normalized benchmark records across 15,526 models from 7 public datasets, we discover a \textit{measurement gap}: 9 of 9 attempts to compute hand-designed phi functions from public benchmarks fail due to metric incomparability, domain transfer breakdown, and dimensional collapse. We then show the gap is specific to interpretable phi functions. Replacing the 16 hand-designed phis with 16 learned bilinear forms trained on 740,803 pairwise routing samples from 5 benchmark sources, we achieve 83.63\% accuracy on RouterBench (across 11 models) and AUC 0.8006 on RouteLLM (retaining 94.35\% of always-strong quality at 50\% cost). A leave-one-out generalization analysis reveals that RouterBench is the critical training source: removing it collapses cross-source transfer, while removing MT-Bench or RewardBench has no measurable effect. A scaling ablation (BilinearPhiEngine V3, 14.6M parameters, 4.63$\times$ V1) trained on an expanded dataset of 1.98 million samples with progressive architectural additions (task conditioning, cross-attention, metadata features) underperforms V1 on every metric, confirming that data quality dominates model capacity for learned routing. The measurement gap thus separates interpretable from learned routing, not theory from practice. Our implementation routes across heterogeneous endpoints (LLMs, agents, scripts, tool-use systems) with literature-informed seed weights for the interpretable interface and bilinear heads for empirical validation.
\end{abstract}

\vspace{6pt}
\hrule
\vspace{12pt}

% === Body ===
\section{Introduction}
\label{sec:intro}

Language model routing, the problem of selecting the optimal model for a given task, has become a practical necessity. With dozens of models available at different price points, quality levels, and capability profiles, manual selection is intractable. Automated routing promises to reduce costs while maintaining quality by matching tasks to appropriate endpoints.

Existing approaches frame routing as either binary classification (strong vs. weak model) \citep{ong2024routellm}, cascading (try cheap first, escalate if needed) \citep{chen2023frugalgpt}, or learned scoring \citep{hu2024routerbench}. These methods share three limitations: they route only between LLMs, ignoring agents, scripts, and tool-use systems; they use additive scoring functions vulnerable to compensating failures; and they lack formal convergence guarantees.

We present S(M,T), a unified scoring framework that addresses all three:

\begin{equation*}
\smt = \underbrace{\left[\prod_{j} g_j(M,T)\right]}_{\text{Gates}} \cdot \underbrace{\left[\sum_{i} w_i \cdot \phii(M,T)\right]}_{\text{Compatibility}} \cdot \underbrace{e^{-\betavec(T) \cdot \mathbf{H}(M)}}_{\text{Cost penalty}}
\end{equation*}

The equation combines three multiplicative components: a Product-of-Experts gate layer that enforces hard constraints (Section~\ref{sec:framework}), a weighted sum of 16 phi compatibility functions capturing different dimensions of model-task fit, and a Boltzmann cost penalty with task-conditional temperature. The multiplicative structure is not decorative: it provably eliminates the equal-scoring vulnerability of additive approaches (Proposition~\ref{prop:poe}).

\paragraph{Theoretical Contributions.}
We prove convergence under Hajek conditions with logarithmic cooling (Theorem~\ref{thm:hajek}), establish $O(\sqrt{KN\log N})$ regret bounds through contextual bandits reduction (Theorem~\ref{thm:regret}), and demonstrate weight and temperature identifiability under standard rank conditions (Theorems~\ref{thm:weight_id},~\ref{thm:beta_id}). Six orthogonality proofs verify that the 16 phi functions capture non-redundant information (Theorem~\ref{thm:orthogonality}). All 14 results verify at 85--95\% confidence.

\paragraph{The Measurement Gap.}
A key empirical finding is negative. We ingested 2,764,077 normalized benchmark records across 15,526 models from 7 public datasets (Section~\ref{sec:empirical}). We then attempted 14 measurement designs to ground hand-designed phi functions in this data (Section~\ref{sec:measurement}). All 14 failed due to structural problems: metric scale incomparability, domain transfer breakdown, benchmark contamination, and dimensional collapse.

\paragraph{Learned Phi Functions.}
However, the measurement gap is specific to interpretable, hand-designed phi functions. We show that replacing the 16 hand-designed phis with 16 learned bilinear forms, trained end-to-end on 740,803 pairwise routing samples, bypasses the measurement gap entirely (Section~\ref{sec:learned_phi}). The BilinearPhiEngine achieves 83.63\% accuracy on RouterBench and AUC 0.8006 on RouteLLM, demonstrating that the weighted-phi architecture produces competitive routing when the learning problem is well-posed. A leave-one-out generalization analysis identifies which training sources matter (RouterBench is critical; MT-Bench and RewardBench are redundant) and which transfer problems remain hard (RouteLLM's binary format resists prediction from other sources). The tradeoff is interpretability: the 16 learned heads carry no semantic labels.

\paragraph{Implementation.}
The v0.1 router implements S(M,T) across 13 heterogeneous endpoints: 7 LLMs, 3 autonomous agents, 1 script, and 2 MCP tool-use systems (Section~\ref{sec:empirical}). The hand-designed phi functions with literature-informed seed weights serve as the interpretable production interface, while the bilinear engine provides empirical validation. Robbins-Monro online learning adapts weights from routing feedback.

\paragraph{Paper Organization.}
Section~\ref{sec:background} reviews background on routing, stochastic approximation, and contextual bandits. Section~\ref{sec:framework} defines the S(M,T) equation, gates, phi functions, and seed weights. Section~\ref{sec:convergence} presents convergence, regret, identifiability, and orthogonality results. Section~\ref{sec:adversarial} analyzes adversarial robustness. Section~\ref{sec:measurement} documents the measurement gap. Section~\ref{sec:empirical} describes the data pipeline and implementation. Section~\ref{sec:learned_phi} presents the learned BilinearPhiEngine with benchmark evaluation and generalization analysis. Section~\ref{sec:related} surveys related work. Section~\ref{sec:discussion} discusses limitations and future directions. Section~\ref{sec:conclusion} concludes.

\section{Background}
\label{sec:background}

\subsection{LLM Routing}

The proliferation of language models with different cost-quality tradeoffs has created a routing problem: given a user query, which model should handle it? Early work framed this as binary classification (strong model vs. weak model), while recent approaches treat it as a multi-armed bandit or learned scoring problem.

\textbf{FrugalGPT} \citep{chen2023frugalgpt} introduced cascading: try a cheap model first, escalate to an expensive model only if confidence is low. This reduces cost by 50--90\% with minimal quality loss on standard benchmarks. However, the cascade structure is rigid: it assumes a fixed ordering of models by cost, and the escalation threshold must be tuned per task type.

\textbf{RouteLLM} \citep{ong2024routellm} trained a classifier to predict whether GPT-4 or a smaller model would better handle a given query. The approach achieves strong results on the specific model pair it was trained for but does not generalize to arbitrary endpoint sets or non-LLM endpoints.

\textbf{RouterBench} \citep{hu2024routerbench} provided a benchmark for evaluating routing algorithms, standardizing the comparison across 11 methods. Their analysis revealed that no single routing method dominates across all task types, motivating multi-dimensional scoring approaches.

\subsection{Scoring Functions for Model Selection}

Existing routing scores are typically additive: $S = \sum_i w_i f_i$ where $f_i$ are features and $w_i$ are weights. This structure has a known vulnerability: a model that catastrophically fails one requirement can still achieve a high total score by excelling at others. The Product-of-Experts framework \citep{hinton2002poe} addresses this through multiplicative combination, where any single expert can veto the prediction.

\subsection{Stochastic Approximation}

The Robbins-Monro algorithm \citep{robbins1951stochastic} provides a framework for online parameter estimation under noisy observations. Given a target function $f(\theta) = 0$, the update $\theta_{n+1} = \theta_n + \alpha_n \cdot (Y_n - f(\theta_n))$ converges to the root under standard conditions (bounded noise, decaying step sizes). This is the theoretical foundation for the weight learning procedure in Section~\ref{sec:convergence}.

\subsection{Simulated Annealing}

\citet{hajek1988cooling} established necessary and sufficient conditions for simulated annealing to converge to the global optimum. The key result: a logarithmic cooling schedule $\beta(t) = c \log(1+t)$ with $c$ exceeding the maximum energy barrier suffices for convergence. We apply this to the temperature schedule in the S(M,T) scoring equation.

\subsection{Contextual Bandits}

In contextual bandits \citep{agrawal2013thompson}, an agent observes a context (task features), selects an action (endpoint), and receives a reward (quality signal). Thompson Sampling with linear payoffs achieves $O(\sqrt{dT\log T})$ regret where $d$ is the context dimension and $T$ is the horizon. This provides the regret bound for our routing framework (Section~\ref{sec:convergence}).

\subsection{Entropy Regularization}

Soft Actor-Critic \citep{haarnoja2018sac} introduced entropy-regularized policy optimization in reinforcement learning, adding $\lambda H(\pi)$ to the objective to prevent premature convergence. The temperature $\lambda$ can be fixed or learned. We adapt this to routing, using a decaying $\lambda(n)$ schedule that transitions from exploration to exploitation.

\section{The S(M,T) Framework}
\label{sec:framework}

We define a unified scoring function that maps any (model, task) pair to a real-valued compatibility score. The framework generalizes beyond LLM-to-LLM routing: $M$ represents any capability endpoint (language model, autonomous agent, script, or tool-use system), and $T$ represents a structured task description.

\subsection{Definitions}

\begin{definition}[Endpoint]
An endpoint $M \in \mathcal{M}$ is a callable capability with associated metadata: context window $c_M$, cost vector $\mathbf{h}_M \in \R^d_+$, capability flags $\mathbf{f}_M \in \{0,1\}^k$, and type $\tau_M \in \{\text{LLM}, \text{Agent}, \text{Script}, \text{Tool}\}$.
\end{definition}

\begin{definition}[Task]
A task $T \in \mathcal{T}$ is a user request with classification $t_T \in \{$general, coding, math, reasoning, knowledge, instruction\_following$\}$, estimated complexity $\kappa_T \in [0,1]$, context length $\ell_T \in \mathbb{N}$, and constraint flags $\mathbf{r}_T \in \{0,1\}^k$.
\end{definition}

\subsection{The Scoring Equation}

\begin{definition}[S(M,T) Score]
For an endpoint $M$ and task $T$, the compatibility score is:
\begin{equation}
\label{eq:smt}
\smt = \underbrace{\left[\prod_{j=1}^{J} g_j(M,T)\right]}_{\text{Gate layer}} \cdot \underbrace{\left[\sum_{i=1}^{K} w_i \cdot \phii(M,T)\right]}_{\text{Weighted compatibility}} \cdot \underbrace{e^{-\betavec(T) \cdot \mathbf{H}(M)}}_{\text{Cost penalty}}
\end{equation}
where $J=5$ binary gates, $K=16$ phi compatibility functions, and $d$-dimensional cost vector.
\end{definition}

The three multiplicative components serve distinct roles:

\paragraph{Gate Layer (Product-of-Experts).} Each gate $g_j: \mathcal{M} \times \mathcal{T} \to \{0, 1\}$ enforces a hard constraint. If any gate returns 0, the entire score is zeroed out. This implements a Product-of-Experts (PoE) architecture \citep{hinton1999poe} that eliminates the equal-scoring vulnerability of purely additive approaches, where a model catastrophically failing one requirement could still score well by excelling at others.

The five gates are:
\begin{enumerate}
    \item \textbf{Context window}: $g_1 = \mathbf{1}[c_M \geq \ell_T]$
    \item \textbf{Format support}: $g_2 = \mathbf{1}[\text{structured output supported if required}]$
    \item \textbf{Availability}: $g_3 = \mathbf{1}[M \text{ is currently reachable}]$
    \item \textbf{Budget}: $g_4 = \mathbf{1}[\text{estimated cost} \leq b_T]$
    \item \textbf{Thinking support}: $g_5 = \mathbf{1}[\text{CoT supported if required}]$
\end{enumerate}

\paragraph{Weighted Compatibility Functions.} The phi functions $\phii: \mathcal{M} \times \mathcal{T} \to [0,1]$ capture different dimensions of model-task compatibility. The weight vector $\wvec = (w_1, \ldots, w_K)$ lies on the probability simplex $\Delta^{K-1}$, so $\sum_i w_i = 1$ and $w_i \geq 0$. This normalization ensures the weighted sum produces a $[0,1]$-bounded compatibility score.

\paragraph{Boltzmann Cost Penalty.} The exponential term $e^{-\betavec(T) \cdot \mathbf{H}(M)}$ penalizes expensive endpoints with a task-conditional temperature. Unlike a fixed temperature, the vector $\betavec(T) \in \R^d_+$ allows different task types to tolerate different cost levels: a safety-critical medical query should accept higher costs than a simple formatting task.

\subsection{The 16 Phi Functions}

Table~\ref{tab:phi_functions} defines all 16 compatibility functions. They decompose into four categories based on data availability and computation type:

\begin{table}[h]
\centering
\caption{Phi function definitions with seed weights and data coverage.}
\label{tab:phi_functions}
\small
\begin{tabular}{llccrr}
\toprule
\textbf{ID} & \textbf{Name} & \textbf{Type} & \textbf{$w_i$ (seed)} & \textbf{Records} & \textbf{Models} \\
\midrule
$\phi_1$ & constraint\_filter & Computed & 0.0706 & 4,200 & 350 \\
$\phi_2$ & domain\_classifier & Data & 0.0353 & 2,042,817 & 14,926 \\
$\phi_3$ & complexity\_estimator & Data & 0.1058 & 2,042,817 & 14,926 \\
$\phi_4$ & performance\_predictor & Data & 0.1764 & 2,759,877 & 15,176 \\
$\phi_5$ & context\_fitness & Prior & 0.0081 & 0 & 0 \\
$\phi_6$ & historical\_success & Data & 0.1411 & 2,307,010 & 14,928 \\
$\phi_7$ & uncertainty\_score & Derived & 0.0163 & 0 & 0 \\
$\phi_8$ & cascade\_position & Data & 0.1235 & 2,311,210 & 15,278 \\
$\phi_9$ & load\_availability & Runtime & 0.0081 & 0 & 0 \\
$\phi_{10}$ & cross\_node\_capability & Prior & 0.0081 & 0 & 0 \\
$\phi_{11}$ & calibration\_quality & Prior & 0.0081 & 0 & 0 \\
$\phi_{12}$ & output\_structure & Prior & 0.0163 & 0 & 0 \\
$\phi_{13}$ & reasoning\_depth & Data & 0.1058 & 2,003,949 & 12,434 \\
$\phi_{14}$ & cost\_efficiency\_ratio & Data & 0.0882 & 719,702 & 576 \\
$\phi_{15}$ & context\_window\_util & Data & 0.0176 & 1,400 & 350 \\
$\phi_{16}$ & instruction\_adherence & Data & 0.0706 & 113,853 & 5,795 \\
\midrule
\multicolumn{4}{l}{\textbf{Total}} & \textbf{2,764,077} & \textbf{15,526} \\
\bottomrule
\end{tabular}
\end{table}

\textbf{Data-grounded} ($\phi_2, \phi_3, \phi_4, \phi_6, \phi_8, \phi_{13}, \phi_{14}, \phi_{16}$): Query 2.76M normalized benchmark records to compute model-specific scores. These carry the highest weight in the seed configuration ($\sum w_i = 0.85$ for data-grounded phis).

\textbf{Computed} ($\phi_1, \phi_{12}, \phi_{15}$): Deterministic functions of endpoint metadata and task properties (e.g., context window ratio, cost tiers).

\textbf{Prior-based} ($\phi_5, \phi_{10}, \phi_{11}$): Return literature-informed default values due to absent public data. These receive low seed weights ($w_i = 0.0081$) reflecting low confidence.

\textbf{Runtime/Derived} ($\phi_7, \phi_9$): Computed at inference time from system state (availability, data coverage).

\subsection{Seed Weight Derivation}

The initial weight vector $\wvec_0$ is not arbitrary. We derive it from two signals:

\begin{equation}
\label{eq:seed_weights}
w_i^{(0)} = \frac{p_i \cdot q_i}{\sum_{k=1}^K p_k \cdot q_k}
\end{equation}

where $p_i$ is the literature prior (importance weight from FrugalGPT, RouteLLM, and RouterBench) on a 0.5--5.0 scale, and $q_i$ is a data quality multiplier:
\[
q_i = \begin{cases}
1.0 & \text{if data quality = high (>1000 records, >50 models)} \\
0.5 & \text{if data quality = medium} \\
0.25 & \text{if data quality = low} \\
0.1 & \text{if data quality = none}
\end{cases}
\]

This places the most weight on phi functions that are both important in the literature \emph{and} well-supported by data. The resulting top-4 seed weights are: $\phi_4$ (performance, $w=0.176$), $\phi_6$ (history, $w=0.141$), $\phi_8$ (cascade, $w=0.124$), and $\phi_3$/$\phi_{13}$ (complexity/reasoning, $w=0.106$ each).

\subsection{Entropy-Regularized Selection}

To prevent premature convergence to a single endpoint (mode collapse), we add an entropy bonus to the routing policy:

\begin{equation}
\label{eq:entropy}
\smt_{\text{reg}} = \smt + \lambda(n) \cdot H(\pi)
\end{equation}

where $H(\pi) = -\sum_M \pi(M) \log \pi(M)$ is the Shannon entropy of the selection distribution, and $\lambda(n)$ decays as:

\begin{equation}
\lambda(n) = \max\left(\lambda_{\min}, \lambda_0 \cdot \gamma^n\right)
\end{equation}

with $\lambda_0 = 0.5$, $\gamma = 0.995$, and $\lambda_{\min} = 0.01$. This schedule, inspired by Soft Actor-Critic \citep{haarnoja2018sac}, maintains exploration early (high $\lambda$) and converges to greedy selection as confidence grows (low $\lambda$).

\subsection{Phi Function Orthogonality}

A natural concern is whether 16 phi functions are redundant. We verify orthogonality through six formal proofs (Section~\ref{sec:convergence}), showing that each phi captures a distinct signal:

\begin{itemize}
    \item $\phi_{11}$ (calibration) $\perp$ $\phi_7$ (uncertainty): Model-internal confidence vs. router-external data sufficiency.
    \item $\phi_{12}$ (structure) $\perp$ $\phi_1$ (constraint): Continuous reliability gradient vs. binary capability check.
    \item $\phi_{13}$ (reasoning) $\perp$ $\phi_3$ (complexity): Multi-step depth vs. surface-level difficulty.
    \item $\phi_{14}$ (cost efficiency) $\perp$ $\beta \cdot H$: Task-conditional quality-per-dollar vs. global cost penalty.
    \item $\phi_{15}$ (context util) $\perp$ $\phi_1$: Performance degradation curve vs. binary fit.
    \item $\phi_{16}$ (instruction) $\perp$ $\phi_5$ (context fitness): Behavioral compliance vs. content relevance.
\end{itemize}

This ensures the 16 dimensions are not collapsible to a smaller set without information loss.

\section{Convergence and Optimality}
\label{sec:convergence}

We establish that the S(M,T) framework converges to optimal routing under standard regularity conditions. This section presents eight verified results: global convergence via simulated annealing (Theorem~\ref{thm:hajek}), sublinear regret (Theorem~\ref{thm:regret}), parameter identifiability (Theorems~\ref{thm:weight_id} and~\ref{thm:beta_id}), the gate layer fix (Proposition~\ref{prop:poe}), informed prior acceleration (Proposition~\ref{prop:prior}), entropy schedule correctness (Proposition~\ref{prop:entropy}), and phi function orthogonality (Theorem~\ref{thm:orthogonality}). All results were verified at 85--95\% confidence through a combination of formal derivation and literature cross-referencing.

\subsection{Convergence under Hajek Conditions}

The weight learning procedure uses Robbins-Monro stochastic approximation \citep{robbins1951stochastic} with a logarithmic cooling schedule.

\begin{definition}[Weight Update Rule]
After observing feedback $r_n$ (quality score) from decision $n$, the weight vector updates as:
\begin{equation}
\label{eq:rm_update}
w_i^{(n+1)} = \Pi_{\Delta^{K-1}}\left[w_i^{(n)} + \alpha_n \cdot \phi_i(M_n, T_n) \cdot (r_n - \hat{r}_n)\right]
\end{equation}
where $\alpha_n = 1/(n+1)$ is the step size, $\hat{r}_n = \sum_i w_i^{(n)} \phi_i(M_n, T_n)$ is the predicted score, and $\Pi_{\Delta^{K-1}}$ projects onto the probability simplex via the algorithm of \citet{duchi2008simplex}.
\end{definition}

\begin{theorem}[Global Convergence]
\label{thm:hajek}
Let the cooling schedule be $\beta(t) = c \cdot \log(1 + t)$ with $c > 0$. Under the following conditions:
\begin{enumerate}
    \item All phi functions are bounded: $\phi_i(M,T) \in [0,1]$ for all $i, M, T$.
    \item Step sizes satisfy the Robbins-Monro conditions: $\sum_{n=1}^{\infty} \alpha_n = \infty$ and $\sum_{n=1}^{\infty} \alpha_n^2 < \infty$.
    \item The endpoint set $\mathcal{M}$ is finite.
\end{enumerate}
Then the weight vector $\wvec^{(n)}$ converges almost surely to the optimal weights $\wvec^*$ that maximize expected routing quality:
\[
\wvec^{(n)} \xrightarrow{a.s.} \wvec^* = \arg\max_{\wvec \in \Delta^{K-1}} \E\left[\text{quality}(\arg\max_M \smt)\right]
\]
\end{theorem}

\begin{proof}[Proof sketch]
The step size $\alpha_n = 1/(n+1)$ satisfies both Robbins-Monro conditions: $\sum 1/(n+1) = \infty$ and $\sum 1/(n+1)^2 = \pi^2/6 < \infty$. Since $\phi_i \in [0,1]$ and the error term $(r_n - \hat{r}_n)$ is bounded by the boundedness of quality signals, the gradient estimate has bounded variance. The simplex projection $\Pi_{\Delta^{K-1}}$ is a contraction, preserving the descent property. By the Hajek sufficient conditions for convergence of simulated annealing \citep{hajek1988cooling}, logarithmic cooling $\beta(t) = c \log(1+t)$ guarantees convergence to the global optimum of the score landscape, provided the energy barriers are finite (which follows from boundedness of phi functions). The Robbins-Monro conditions then ensure the stochastic gradient noise averages out, yielding almost sure convergence by the ODE method of \citet{kushner2003stochastic}.
\end{proof}

\subsection{Regret Bounds via Contextual Bandits}

The routing problem maps naturally to contextual bandits: at each round, the router observes a context (task $T$), selects an arm (endpoint $M$), and receives a reward (quality signal $r$). The phi function values $(\phi_1(M,T), \ldots, \phi_K(M,T))$ form the feature vector for the linear payoff model.

\begin{theorem}[Sublinear Regret]
\label{thm:regret}
Let $|\mathcal{M}| = L$ endpoints and $K$ phi functions. Under Thompson Sampling with linear payoffs \citep{agrawal2013thompson}, the cumulative regret after $N$ routing decisions satisfies:
\begin{equation}
\label{eq:regret}
R(N) = O\left(\sqrt{KN\log N}\right)
\end{equation}
For the S(M,T) configuration with $K=16$ and $L=13$, achieving 95\% optimal routing requires approximately $N \approx 10{,}000$ decisions.
\end{theorem}

\begin{proof}[Proof sketch]
This follows from the regret analysis of Thompson Sampling for contextual bandits with linear payoffs \citep{agrawal2013thompson}. The key insight is that S(M,T) is linear in the weights $\wvec$ for fixed phi values. Specifically, the expected reward for selecting endpoint $M$ on task $T$ is $\E[r | M, T] = \wvec^{\top} \boldsymbol{\phi}(M,T) + \epsilon$ where $\boldsymbol{\phi}(M,T) = (\phi_1(M,T), \ldots, \phi_K(M,T))^{\top}$. The $\sqrt{K}$ factor reflects the dimension of the weight space, $\sqrt{N}$ the horizon, and $\log N$ the confidence widths in posterior sampling. The sample complexity estimate follows from setting $R(N)/N \leq 0.05$ and solving for $N$.
\end{proof}

\begin{remark}
The $O(\sqrt{KN\log N})$ bound improves over naive uniform exploration, which requires $O(L \cdot N)$ regret. The contextual structure (using phi functions as features rather than treating each endpoint independently) provides the improvement factor of $\sqrt{K/L}$ when $K < L$.
\end{remark}

\subsection{Parameter Identifiability}

A practical concern: can the weights $\wvec$ and temperature $\betavec$ be uniquely recovered from observed routing outcomes?

\begin{theorem}[Weight Identifiability]
\label{thm:weight_id}
Given $N$ observed score-task pairs $\{(\smt_n, T_n, M_n)\}_{n=1}^N$, the weight vector $\wvec$ is uniquely identifiable if and only if the phi matrix
\[
\Phi = \begin{pmatrix}
\phi_1(M_1, T_1) & \cdots & \phi_K(M_1, T_1) \\
\vdots & \ddots & \vdots \\
\phi_1(M_N, T_N) & \cdots & \phi_K(M_N, T_N)
\end{pmatrix} \in \R^{N \times K}
\]
has full column rank, i.e., $\text{rank}(\Phi) = K$.
\end{theorem}

\begin{proof}[Proof sketch]
Fixing the gates and temperature terms, the observed scores satisfy $\mathbf{s} = \Phi \wvec + \boldsymbol{\epsilon}$, a standard linear regression. Full column rank of $\Phi$ is necessary and sufficient for $(\Phi^{\top}\Phi)^{-1}$ to exist, yielding a unique least-squares solution $\hat{\wvec} = (\Phi^{\top}\Phi)^{-1}\Phi^{\top}\mathbf{s}$ \citep{walter1997identification}. The minimum requirement is $N \geq K = 16$ distinct model-task pairs. With 15,526 unique models in the database, this condition is easily satisfied in practice.
\end{proof}

\begin{theorem}[Temperature Identifiability]
\label{thm:beta_id}
Let $\betavec(T)$ be task-conditional. Given sufficient task diversity (i.e., the task feature matrix has full rank), the temperature parameters are jointly identifiable with the weights.
\end{theorem}

\begin{proof}[Proof sketch]
The score equation $\smt = [\sum_i w_i \phi_i] \cdot e^{-\betavec(T) \cdot \mathbf{H}(M)}$ is separable: the log-score decomposes as $\log \smt = \log(\sum_i w_i \phi_i) - \betavec(T) \cdot \mathbf{H}(M)$. Given diverse tasks with varied cost sensitivities, the Jacobian of the score function with respect to $(\wvec, \betavec)$ has full rank \citep{rothenberg1971identification}. Task-conditional temperature vectors improve identifiability over a scalar $\beta$ because different task types provide independent equations constraining $\betavec$.
\end{proof}

\subsection{Product-of-Experts Gate Layer}

The multiplicative gate structure addresses a fundamental vulnerability of purely additive scoring.

\begin{proposition}[PoE Eliminates Equal-Scoring Vulnerability]
\label{prop:poe}
Consider two endpoints $M_a, M_b$ where $M_a$ fails a hard constraint (e.g., insufficient context window) but excels on all other dimensions, and $M_b$ passes all constraints with moderate scores. Under an additive score $S_{\text{add}} = \sum_i w_i \phi_i$, there exist weight vectors where $S_{\text{add}}(M_a, T) \geq S_{\text{add}}(M_b, T)$. Under the gated score $\smt$, $S(M_a, T) = 0 < S(M_b, T)$ always.
\end{proposition}

\begin{proof}
If $M_a$ fails gate $g_j$, then $g_j(M_a, T) = 0$, so $\prod_j g_j(M_a, T) = 0$, and $S(M_a, T) = 0$ regardless of phi values. Since $M_b$ passes all gates, $\prod_j g_j(M_b, T) = 1$, and $S(M_b, T) = [\sum_i w_i \phi_i(M_b, T)] \cdot e^{-\betavec \cdot \mathbf{H}(M_b)} > 0$ whenever at least one $\phi_i > 0$ and $w_i > 0$.

For the additive case, let $M_a$ have $\phi_i(M_a, T) = 1$ for the top-weighted functions. Since weights sum to 1, $S_{\text{add}}(M_a, T)$ can approach 1.0 while $S_{\text{add}}(M_b, T)$ with moderate phi values stays below. The gated formulation eliminates this failure mode entirely. This is the Product-of-Experts principle: any single expert (gate) can veto \citep{hinton2002poe}.
\end{proof}

\begin{remark}
The gate layer preserves differentiability in the valid region. For all endpoints where every gate returns 1, the score function is continuously differentiable in $\wvec$ and $\betavec$, so gradient-based learning remains well-defined.
\end{remark}

\subsection{Informed Prior Acceleration}

\begin{proposition}[Prior-Informed Convergence Speed]
\label{prop:prior}
Let $\wvec_0$ be the seed weight vector derived from literature priors and data quality multipliers (Equation~\ref{eq:seed_weights}), and let $\wvec_U = (1/K, \ldots, 1/K)$ be the uniform initialization. Under the Robbins-Monro update (Equation~\ref{eq:rm_update}), the expected regret with informed prior satisfies:
\begin{equation}
\E[R_{\text{informed}}(N)] \leq \E[R_{\text{uniform}}(N)] - \Omega\left(\|\wvec_U - \wvec^*\|^2 - \|\wvec_0 - \wvec^*\|^2\right)
\end{equation}
provided $\|\wvec_0 - \wvec^*\| < \|\wvec_U - \wvec^*\|$.
\end{proposition}

\begin{proof}[Proof sketch]
The cumulative regret of Robbins-Monro is dominated by the initial transient, which scales as $\|\wvec^{(0)} - \wvec^*\|^2 / \alpha_1$ \citep{auer2002finite}. Since both initializations use the same step size schedule, the regret difference is proportional to the squared distance difference. The condition $\|\wvec_0 - \wvec^*\| < \|\wvec_U - \wvec^*\|$ holds when the literature priors are positively correlated with the true optimal weights, which is the standard assumption for informed Bayesian initialization.
\end{proof}

The seed weight vector $\wvec_0$ places 85\% of total weight on data-grounded phi functions ($\phi_2, \phi_3, \phi_4, \phi_6, \phi_8, \phi_{13}, \phi_{14}, \phi_{16}$), matching the empirical finding that benchmark-informed functions carry the strongest signal. The top four seed weights ($\phi_4 = 0.176$, $\phi_6 = 0.141$, $\phi_8 = 0.124$, $\phi_3 = \phi_{13} = 0.106$) align with the literature consensus that performance prediction and historical success are the primary routing signals \citep{chen2023frugalgpt, ong2024routellm}.

\subsection{Entropy Regularization Schedule}

\begin{proposition}[Entropy Schedule Correctness]
\label{prop:entropy}
The decaying entropy regularization $\lambda(n) = \max(\lambda_{\min}, \lambda_0 \cdot \gamma^n)$ with $\lambda_0 = 0.5$, $\gamma = 0.995$, and $\lambda_{\min} = 0.01$ satisfies:
\begin{enumerate}
    \item \textbf{Early exploration}: For $n < 100$, $\lambda(n) > 0.3$, ensuring the entropy bonus dominates score differences smaller than $0.3 \cdot \log |\mathcal{M}|$.
    \item \textbf{Asymptotic exploitation}: $\lambda(n) \to \lambda_{\min} = 0.01$, so the regularized score $\smt_{\text{reg}} \approx \smt$ for large $n$.
    \item \textbf{Convergence preservation}: Adding $\lambda(n) H(\pi)$ does not affect the convergence guarantee of Theorem~\ref{thm:hajek}, since $\lambda(n) \to \lambda_{\min}$ and the entropy term is bounded.
\end{enumerate}
\end{proposition}

\begin{proof}
(1) At $n=100$: $\lambda(100) = 0.5 \cdot 0.995^{100} = 0.5 \cdot 0.606 = 0.303 > 0.3$.

(2) The decay $\gamma^n \to 0$ as $n \to \infty$, so $\lambda(n) \to \lambda_{\min} = 0.01$. With $|\mathcal{M}| = 13$ endpoints, $\max H(\pi) = \log 13 \approx 2.56$, giving a maximum entropy bonus of $0.01 \cdot 2.56 = 0.026$, which is negligible compared to typical score separations.

(3) The Robbins-Monro convergence in Theorem~\ref{thm:hajek} requires bounded perturbations. Since $\lambda(n) H(\pi) \leq 0.5 \cdot \log 13 < 1.3$ for all $n$, and $\lambda(n) \to 0.01$ geometrically, the perturbation sequence is summable: $\sum_n |\lambda(n) - \lambda_{\min}| < \infty$. By the robustness of stochastic approximation to summable perturbations \citep{kushner2003stochastic}, convergence is preserved.

This schedule mirrors the temperature adaptation in Soft Actor-Critic \citep{haarnoja2018sac}, which uses entropy regularization to balance exploration and exploitation in continuous action spaces.
\end{proof}

\subsection{Phi Function Orthogonality}

We verify that the 16 phi functions capture non-redundant information, establishing that no phi can be expressed as a linear combination of others in the critical pairs.

\begin{theorem}[Pairwise Orthogonality of Critical Pairs]
\label{thm:orthogonality}
The following six pairs of phi functions are orthogonal in the sense that neither can be expressed as a monotone function of the other, and there exist task-model pairs where their rankings diverge:
\end{theorem}

We prove each pair by constructing separating examples.

\paragraph{Proof 1: $\phi_{11}$ (calibration) $\perp$ $\phi_7$ (uncertainty).}
$\phi_{11}$ measures a model's self-knowledge (expected calibration error, ECE), while $\phi_7$ measures the router's epistemic uncertainty about the model. Consider: Model $A$ is well-calibrated ($\phi_{11} = 0.9$) but the router has limited data on it ($\phi_7 = 0.2$). Model $B$ is poorly calibrated ($\phi_{11} = 0.3$) but has extensive routing history ($\phi_7 = 0.9$). The rankings diverge: $\phi_{11}$ prefers $A$, $\phi_7$ prefers $B$. This separation arises because model-internal confidence and router-external data sufficiency are fundamentally different information sources.

\paragraph{Proof 2: $\phi_{12}$ (structure) $\perp$ $\phi_1$ (constraint).}
$\phi_1$ is a binary gate: does the model support the required output format? $\phi_{12}$ is a continuous reliability gradient: given format support, how often does the model produce valid structured output? Two models both supporting JSON mode ($\phi_1 = 1$) can differ sharply: Model $A$ produces invalid JSON 30\% of the time ($\phi_{12} = 0.7$) while Model $B$ achieves 99\% validity ($\phi_{12} = 0.99$). $\phi_1$ cannot distinguish them.

\paragraph{Proof 3: $\phi_{13}$ (reasoning) $\perp$ $\phi_3$ (complexity).}
$\phi_3$ estimates surface-level task difficulty (token count, vocabulary complexity, domain specificity). $\phi_{13}$ measures multi-step reasoning depth (chain length, logical dependency structure). The task ``Prove that the sum of the first $n$ odd numbers equals $n^2$'' has moderate surface complexity ($\phi_3 \approx 0.4$) but requires deep reasoning ($\phi_{13} \approx 0.9$). Conversely, a task requiring specialized domain vocabulary scores high on $\phi_3$ but may need no multi-step reasoning.

\paragraph{Proof 4: $\phi_{14}$ (cost efficiency) $\perp$ $\beta \cdot \mathbf{H}$ (cost penalty).}
The Boltzmann cost penalty $e^{-\betavec \cdot \mathbf{H}}$ is a global, task-agnostic penalty on endpoint cost. $\phi_{14}$ is a task-conditional quality-per-dollar ratio. For a summarization task, a cheap model may achieve 95\% of the quality of an expensive model ($\phi_{14,\text{cheap}} = 0.95$). For a code debugging task, the cheap model achieves only 20\% quality ($\phi_{14,\text{cheap}} = 0.2$). The global penalty $\beta \cdot H$ applies identically to both tasks, while $\phi_{14}$ captures the task-specific efficiency inversion.

\paragraph{Proof 5: $\phi_{15}$ (context utilization) $\perp$ $\phi_1$ (constraint).}
$\phi_1$ is a binary check: does the model's context window fit the input? $\phi_{15}$ captures the continuous degradation of performance as context usage increases. A model with a 128K context window passes $\phi_1$ for a 50K token input, but may suffer 40\% accuracy degradation at that utilization level ($\phi_{15} = 0.6$), while another model with identical context window maintains 95\% accuracy ($\phi_{15} = 0.95$). The binary gate is blind to this quality gradient.

\paragraph{Proof 6: $\phi_{16}$ (instruction adherence) $\perp$ $\phi_5$ (context fitness).}
$\phi_5$ measures content relevance (how well the model's training distribution matches the task domain). $\phi_{16}$ measures behavioral compliance with specific formatting and constraint instructions. Consider the prompt: ``Write a 5-7-5 haiku about machine learning. Do not use the word `AI'.'' A model may produce a beautiful, domain-relevant poem ($\phi_5 = 0.9$) but with wrong syllable counts and containing the forbidden word ($\phi_{16} = 0.2$). $\phi_5$ captures semantic relevance; $\phi_{16}$ captures constraint satisfaction. These are orthogonal properties, and $\phi_{16}$ cannot be expressed as a linear combination of $\phi_3$ (complexity) or $\phi_4$ (performance).

\begin{remark}
These six proofs cover the pairs most susceptible to conflation. The remaining pairs among the 16 phi functions have more obvious separations (e.g., $\phi_9$ (load availability) is a runtime metric independent of all static phi functions). The full $\binom{16}{2} = 120$ pairwise analysis is not necessary: the six critical pairs establish that the phi space is not collapsible to a lower-dimensional representation without information loss.
\end{remark}

\subsection{Summary of Verification Status}

Table~\ref{tab:proofs} summarizes all convergence and optimality results. Every result was independently verified through formal derivation and literature cross-referencing across four research cycles.

\begin{table}[h]
\centering
\caption{Verification status of theoretical results. All 14 results verified at 85--95\% confidence.}
\label{tab:proofs}
\small
\begin{tabular}{llcc}
\toprule
\textbf{Result} & \textbf{Type} & \textbf{Confidence} & \textbf{Key Citation} \\
\midrule
Theorem~\ref{thm:hajek}: Global convergence & Convergence & 90\% & \citet{hajek1988cooling} \\
Theorem~\ref{thm:regret}: $O(\sqrt{KN\log N})$ regret & Regret bound & 90\% & \citet{agrawal2013thompson} \\
Theorem~\ref{thm:weight_id}: Weight identifiability & Identifiability & 85\% & \citet{walter1997identification} \\
Theorem~\ref{thm:beta_id}: Temperature identifiability & Identifiability & 85\% & \citet{rothenberg1971identification} \\
Proposition~\ref{prop:poe}: PoE gate fix & Correctness & 95\% & \citet{hinton2002poe} \\
Proposition~\ref{prop:prior}: Prior acceleration & Convergence & 85\% & \citet{auer2002finite} \\
Proposition~\ref{prop:entropy}: Entropy schedule & Correctness & 85\% & \citet{haarnoja2018sac} \\
Theorem~\ref{thm:orthogonality}: 6 orthogonality proofs & Independence & 90--95\% & Constructive \\
\bottomrule
\end{tabular}
\end{table}

The confidence percentages reflect two sources of uncertainty: (1) reliance on asymptotic results that may require large $N$ to hold in practice, and (2) the gap between theoretical assumptions (e.g., i.i.d. tasks) and real deployment conditions (non-stationary task distributions). Section~\ref{sec:measurement} shows that this gap between theory and measurement, not the theory itself, is the binding constraint.

\section{Adversarial Analysis}
\label{sec:adversarial}

We analyze three attack vectors against the S(M,T) routing framework and show that existing components provide built-in mitigations.

\subsection{Mode Collapse}

\textbf{Attack.} Without regularization, the router converges to always selecting a single endpoint, ignoring potentially better alternatives for specific task types. This is the standard exploitation-exploration failure in bandit settings.

\textbf{Mitigation.} The entropy regularization term (Equation~\ref{eq:entropy}) directly addresses mode collapse. With $\lambda_0 = 0.5$, the entropy bonus during early routing is comparable in magnitude to score differences between endpoints. For 13 endpoints, $\max H(\pi) = \log 13 \approx 2.56$, so the initial entropy bonus can reach $0.5 \times 2.56 = 1.28$, which dominates typical score separations of 0.1--0.3.

As $\lambda(n)$ decays toward $\lambda_{\min} = 0.01$, the router gradually transitions from exploration to exploitation. Proposition~\ref{prop:entropy} establishes that this decay is fast enough to converge but slow enough to prevent premature commitment. The minimum entropy floor $\lambda_{\min} > 0$ ensures the router never fully degenerates to deterministic selection, maintaining a small probability of exploring alternatives indefinitely.

\subsection{Ranking Instability under Perturbation}

\textbf{Attack.} Small changes in phi function outputs cause large changes in endpoint rankings, making the router unreliable. This is particularly concerning when phi functions are noisy (as documented in Section~\ref{sec:measurement}).

\textbf{Mitigation.} The task-conditional temperature vector $\betavec(T)$ provides stability through two mechanisms:

\begin{enumerate}
    \item \textbf{Temperature smoothing.} The Boltzmann cost penalty $e^{-\betavec(T) \cdot \mathbf{H}(M)}$ acts as a soft regularizer. For endpoints with similar phi scores, the temperature term breaks ties based on cost, which is a stable, well-measured quantity (unlike benchmark-derived phi values).

    \item \textbf{Task-conditional sensitivity.} Different task types tolerate different levels of score perturbation. A safety-critical medical query should have a high $\beta$ (strong cost tolerance, selecting the best model regardless of price), while a simple formatting task should have a low $\beta$ (preferring cheaper endpoints). This task-conditional structure means ranking instability in one task type does not propagate to others.
\end{enumerate}

The Product-of-Experts gate layer (Proposition~\ref{prop:poe}) provides additional stability: endpoints that fail hard constraints are eliminated before scoring, removing the most volatile rankings (where a model's score could swing from high to zero depending on whether a gate marginally passes or fails).

\subsection{Weight Manipulation via Feedback}

\textbf{Attack.} An adversary submits systematically biased feedback signals to steer the Robbins-Monro weight updates toward a preferred endpoint. For example, always reporting high quality for a cheap but low-quality model.

\textbf{Mitigation.} Three properties of the learning system limit this attack:

\begin{enumerate}
    \item \textbf{Decaying step size.} The step size $\alpha_n = 1/(n+1)$ means early feedback has outsized influence, but late feedback has diminishing effect. After 1,000 decisions, each new feedback signal shifts weights by at most $0.1\%$. An adversary arriving late faces prohibitive cost to overcome accumulated honest feedback.

    \item \textbf{Simplex constraint.} The Duchi projection onto $\Delta^{K-1}$ prevents any single weight from being driven to extreme values. Inflating one weight necessarily deflates others, limiting the damage from biased feedback on a single phi function.

    \item \textbf{Gate layer independence.} The binary gates are not affected by weight learning. Even if an adversary successfully manipulates weights to favor a particular model, the gate layer still eliminates that model for tasks where it fails hard constraints. The manipulation cannot override context window limits, format support, or budget constraints.
\end{enumerate}

\subsection{Summary}

The three attack vectors map to three defense mechanisms already present in the framework:

\begin{table}[h]
\centering
\caption{Attack vectors and built-in mitigations.}
\label{tab:adversarial}
\small
\begin{tabular}{lll}
\toprule
\textbf{Attack} & \textbf{Defense} & \textbf{Component} \\
\midrule
Mode collapse & Entropy regularization & Equation~\ref{eq:entropy} \\
Ranking instability & Task-conditional $\betavec(T)$ + gates & Equation~\ref{eq:smt} \\
Weight manipulation & Decaying $\alpha_n$ + simplex + gates & Equation~\ref{eq:rm_update} \\
\bottomrule
\end{tabular}
\end{table}

These mitigations are not post-hoc patches. They emerge from the framework's multiplicative structure: gates provide hard safety guarantees that soft scoring cannot override, entropy regularization is a standard technique from reinforcement learning \citep{haarnoja2018sac}, and decaying step sizes are inherent to Robbins-Monro convergence \citep{robbins1951stochastic}.

\section{The Measurement Gap}
\label{sec:measurement}

The preceding sections establish that S(M,T) is theoretically sound: it converges, has sublinear regret, and its parameters are identifiable. This section documents a systematic failure to ground the phi functions in public benchmark data. We attempted 9 measurement designs across 4 research cycles. All failed. The gap between verified theory and measurable practice is the central finding of this paper.

\subsection{Experimental Setup}

We ingested 7 public datasets totaling 2,764,077 normalized benchmark records across 15,526 unique models (Table~\ref{tab:datasets}). Each dataset was downloaded, parsed, and normalized into a unified SQLite schema mapping \texttt{(model\_id, benchmark, metric, score)} tuples. The normalization pipeline applied min-max scaling within each benchmark to produce $[0,1]$-bounded scores.

\begin{table}[h]
\centering
\caption{Public datasets ingested for phi function grounding.}
\label{tab:datasets}
\small
\begin{tabular}{llrrr}
\toprule
\textbf{Dataset} & \textbf{Metric Type} & \textbf{Records} & \textbf{Models} & \textbf{Priority} \\
\midrule
Open LLM Leaderboard & Accuracy (0--100) & 1,897,419 & 5,168 & 3 \\
RouterEval & Pairwise preference & 105,042 & 7,117 & 1 \\
RouteLLM & Win rate & 714,606 & 2 & 2 \\
Epoch AI & Mixed & 40,356 & 2,641 & 4 \\
OpenRouter & Pricing/context & 4,200 & 350 & 5 \\
AlpacaEval & Win rate & 966 & 102 & 11 \\
SWE-bench & Pass rate & 1,488 & 149 & 12 \\
\midrule
\multicolumn{2}{l}{\textbf{Total (deduplicated)}} & \textbf{2,764,077} & \textbf{15,526} & \\
\bottomrule
\end{tabular}
\end{table}

For each phi function, we designed a measurement procedure: a concrete algorithm mapping the normalized database records to a $[0,1]$ score for a given (model, task) pair. Each design was then stress-tested through adversarial analysis, checking for metric conflicts, domain transfer failures, and dimensional errors. A measurement design was \emph{refuted} if we could demonstrate a structural flaw making it unreliable for routing decisions.

\subsection{Nine Measurement Failures}

We present the 9 measurement designs attempted in the final research cycle, all refuted at 85--95\% confidence. Four earlier designs (from the third cycle) were also refuted, giving a cumulative record of 0 verified out of 14 attempted.

\subsubsection{$\phi_2$: Domain Classifier (Refuted, 92\%)}

\textbf{Design.} Classify task $T$ into a domain (code, math, medical, legal, creative, general) using keyword and embedding classifiers. Look up model $M$'s benchmark scores in that domain. Return the normalized score.

\textbf{Failure.} Three irreconcilable problems:
\begin{enumerate}
    \item \textbf{Task-domain ambiguity.} Real tasks are frequently hybrid. A ``medical coding'' task spans both medical and programming domains. Keyword and embedding classifiers cannot resolve multi-domain membership into a single lookup key.
    \item \textbf{Metric scale incomparability.} The domain lookup requires combining scores from benchmarks with fundamentally different scales: MMLU uses accuracy (0--100), Chatbot Arena uses Elo (unbounded, $\sim$1000--2000), and LiveCodeBench uses pass@k (0--1). No statistically sound conversion exists between these scales \citep{deng2023leaderboards}.
    \item \textbf{Normalization destroys semantics.} Min-max normalization within each benchmark produces $[0,1]$ values, but a normalized 0.8 on MMLU and a normalized 0.8 on Chatbot Arena Elo do not represent equivalent capability levels. The normalization creates an illusion of comparability.
\end{enumerate}

\subsubsection{$\phi_{12}$: Output Structure Compliance (Refuted, 92\%)}

\textbf{Design.} Query the Berkeley Function Calling Leaderboard for model $M$'s structured output success rate. Return as $\phi_{12}$.

\textbf{Failure.} Schema complexity degrades reliability non-linearly. A model achieving 95\% success on simple JSON schemas may drop to 40\% on nested schemas with recursive references. The benchmark measures one complexity level; the router needs predictions across all levels. Additionally, keyword-based format detection has catastrophic false negatives: models may produce syntactically valid structured output that pattern matching fails to recognize.

\subsubsection{$\phi_{14}$: Cost Efficiency Ratio (Refuted, 95\%)}

\textbf{Design.} Compute quality-per-dollar as $\phi_{14}(M,T) = \phi_4(M,T) / \text{cost}(M,T)$, where quality comes from benchmark scores and cost from OpenRouter pricing.

\textbf{Failure.} This is the highest-confidence refutation. The ratio is dimensionally invalid: quality $\phi_4$ is a dimensionless benchmark score (0--1), while cost is market-driven (\$/token). Their ratio has no stable physical interpretation. Worse, task-length sensitivity causes $>$100$\times$ cost prediction errors. A legal document analysis (50K tokens) on GPT-4 costs $>$\$50, while a tweet classification (100 tokens) costs $<$\$0.01. Any fixed-length assumption fails catastrophically. Hidden infrastructure costs (70B models require 8-GPU servers) add further multiplicative errors not captured by per-token pricing.

\subsubsection{$\phi_{16}$: Instruction Adherence (Refuted, 92\%)}

\textbf{Design.} Average model $M$'s scores across IFEval (verifiable constraints), MT-Bench (open-ended quality), and AlpacaEval (pairwise preference) to produce a single instruction-following score.

\textbf{Failure.} Instruction adherence is not a single dimension. IFEval measures whether specific verifiable constraints are satisfied (word count limits, format requirements). MT-Bench and AlpacaEval measure stylistic quality and coherence. These create conflicting optimization signals: a model may excel at verifiable constraints while producing stilted prose, or write beautifully while ignoring explicit format requirements. Averaging across these benchmarks produces a number that predicts neither behavior reliably. Novel instruction types (legal compliance, domain-specific constraints) have no benchmark coverage at all.

\subsubsection{$\phi_1$ v2: Constraint Filter Redesign (Refuted, 85\%)}

\textbf{Design.} Replace the original token-counting approach with a compositional parser that checks static capability matrices from model documentation (context window, supported formats, available tools).

\textbf{Failure.} Static capability matrices systematically overstate actual capabilities. Model documentation reports maximum context windows, not effective context windows (where quality degrades). Public documentation reflects marketing incentives, not measured performance \citep{liu2023lost}. Compositional constraints (``output JSON with nested arrays where each element follows schema X'') require understanding the model's reasoning capabilities, which documentation does not capture.

\subsubsection{$\phi_5$ v2: Context Fitness Redesign (Refuted, 92\%)}

\textbf{Design.} Replace embedding cosine similarity with a topic classifier: classify task $T$ by topic, look up model $M$'s historical success rate on that topic from benchmark data.

\textbf{Failure.} Benchmark scores measure broad domain knowledge, not task-specific fitness. A model scoring well on MMLU-medical does not necessarily handle a specific clinical reasoning task well. Static benchmarks also suffer from data contamination: modern LLMs are trained on benchmark datasets, inflating scores beyond true capability \citep{deng2023leaderboards}. Domain classification errors from $\phi_2$ propagate catastrophically into $\phi_5$, compounding the measurement failure.

\subsubsection{$\phi_{13}$ v2: Reasoning Depth Redesign (Refuted, 85\%)}

\textbf{Design.} Analyze GSM8K solution step counts per model. Use step count as a proxy for multi-step reasoning capability. Generalize to other task types.

\textbf{Failure.} Surface features (mathematical operator count, question length) correlate poorly with actual reasoning depth. The task ``prove $\sum_{k=1}^n (2k-1) = n^2$'' has few operators but requires inductive reasoning. GSM8K performance does not generalize to non-mathematical reasoning: legal argument construction, scientific hypothesis testing, and strategic planning require qualitatively different reasoning structures. Normalization across GSM8K, MATH, and BBH introduces further artifacts from incompatible evaluation scales.

\subsubsection{$\phi_{15}$ v2: Context Window Utilization Redesign (Refuted, 92\%)}

\textbf{Design.} Fit degradation curves from RULER benchmark data. For a given model $M$ and context length $\ell_T$, interpolate the expected performance from the fitted curve.

\textbf{Failure.} Task-type performance variance dominates context-length effects. RULER shows $>$30\% performance gaps between retrieval and question-answering tasks at identical context lengths. A single degradation curve per model averages over these gaps, destroying the task-conditional signal. Positional sensitivity compounds the problem: the ``Lost in the Middle'' phenomenon \citep{liu2023lost} shows up to 60\% performance drops when critical information sits in the context middle versus the ends. The available data (5--7 context length measurement points per model) is insufficient to fit the resulting non-monotonic, task-dependent, position-sensitive degradation surface.

\subsubsection{$\phi_4$: Performance Predictor (Inconclusive, 30\%)}

\textbf{Design.} Use the full normalized benchmark database (2.76M records) to predict model $M$'s performance on task type $t_T$.

\textbf{Failure.} Verification failed due to technical issues (JSON parse errors in the automated verification pipeline). However, the same structural problems affecting all other measurement designs apply: benchmark-to-task domain transfer breakdown, metric normalization artifacts, and data contamination. This phi function has the richest data backing (15,176 models) yet could not be verified, suggesting the problem is fundamental rather than data-limited.

\subsection{Root Cause Analysis}

The 9 failures share five structural root causes that are unlikely to be resolved by collecting more data or improving normalization algorithms.

\paragraph{1. Metric Scale Incomparability.} Public benchmarks use fundamentally different measurement scales: accuracy percentages, Elo ratings, pass@k rates, and pairwise win rates. No statistically sound conversion exists between these scales. Min-max normalization creates comparable numbers but destroys the semantic content that would make them useful for routing.

\paragraph{2. Domain Transfer Breakdown.} Performance on curated benchmark tasks does not predict performance on real user tasks. Benchmarks are designed for model comparison, not task-specific prediction. The evaluation conditions (fixed prompts, controlled difficulty, clean data) differ systematically from deployment conditions (varied prompts, uncalibrated difficulty, noisy context).

\paragraph{3. Dimensional Collapse.} Complex, multi-dimensional phenomena (reasoning depth, instruction adherence, cost efficiency) cannot be reduced to single scalar scores without losing the task-conditional structure that routing requires. A model's reasoning capability is not a number; it is a function of task type, problem structure, and prompting strategy.

\paragraph{4. Benchmark Contamination.} Modern LLMs are trained on datasets that include benchmark test sets, inflating scores beyond true capability \citep{deng2023leaderboards}. This contamination is not uniformly distributed: some models are more contaminated than others, some benchmarks more leaked than others. The resulting scores carry systematic bias that normalization cannot remove.

\paragraph{5. Sparse Measurement vs. Complex Phenomena.} The available benchmark data provides sparse coverage (5--7 measurement points for context degradation, single accuracy numbers for domain performance) of phenomena that require dense, multi-dimensional characterization. Fitting complex surfaces to sparse data produces unreliable interpolation at exactly the operating points routing decisions depend on.

\subsection{The Gap is the Finding}

The measurement gap is a central empirical finding, but it requires careful scoping. The gap applies specifically to \emph{hand-designed, interpretable} phi functions that require explicit metric normalization, domain classification, and semantic mapping.

\begin{itemize}
    \item \textbf{Theory}: 14 of 14 mathematical results verified at 85--95\% confidence. The S(M,T) equation converges, has bounded regret, and its parameters are identifiable.
    \item \textbf{Hand-designed measurement}: 0 of 14 measurement designs verified. Every attempt to compute interpretable phi functions from public benchmark data failed due to structural (not data-quantity) limitations.
    \item \textbf{Learned measurement}: The BilinearPhiEngine (Section~\ref{sec:learned_phi}) achieves 83.63\% routing accuracy on RouterBench and AUC 0.8006 on RouteLLM by learning phi functions end-to-end from the \emph{same} benchmark data. The gap is not in the data. It is in the requirement that phi functions be interpretable.
\end{itemize}

This distinction reframes the finding. The binding constraint on LLM routing quality is not the routing algorithm, and it is not the data. It is the demand for interpretable, semantically labeled compatibility functions. When that demand is relaxed, public benchmark data is sufficient. When it is maintained, the five root causes above make static measurement intractable.

\subsection{Implications for the Field}

The measurement gap, and its learned resolution, have four consequences for LLM routing research:

\paragraph{Public benchmarks are insufficient for interpretable routing.} The 7 datasets we ingested represent the most comprehensive public collection available for routing research. If 2.76M records across 15,526 models cannot ground even one hand-designed phi function, the problem is structural for interpretable approaches. Routing-specific benchmarks, ones that measure task-conditional, model-specific performance under deployment conditions, would help.

\paragraph{Learned routing functions bypass the gap.} Section~\ref{sec:learned_phi} demonstrates that bilinear forms trained on pairwise routing data achieve competitive accuracy without requiring explicit metric normalization or domain classification. The normalization and domain transfer problems disappear because the model learns whatever internal representation best predicts routing outcomes. This suggests learned approaches may be more practical than hand-designed ones for production routing.

\paragraph{Online learning is necessary, not optional.} Since static benchmark data cannot reliably initialize interpretable phi functions, the router must learn from its own routing decisions. The Robbins-Monro framework (Section~\ref{sec:convergence}) with entropy regularization provides the theoretical foundation for this learning. The seed weights serve as an informed prior, not as ground truth.

\paragraph{Transparency about measurement quality matters.} Routing papers that report benchmark-grounded scores without documenting the normalization decisions, domain transfer assumptions, and contamination risks are overstating their confidence. The confidence column in Table~\ref{tab:phi_functions} (Section~\ref{sec:framework}), distinguishing data-grounded from prior-based phi functions, is not a weakness of our framework but a design feature.

\section{Empirical Grounding}
\label{sec:empirical}

We describe the data infrastructure supporting the S(M,T) framework: the ingestion pipeline, normalized database, and seed weight computation. While Section~\ref{sec:measurement} establishes that public benchmarks cannot reliably ground individual phi functions, the aggregate data still informs the seed weight vector $\wvec_0$ and provides the initial endpoint rankings used during the cold-start period.

\subsection{Data Ingestion Pipeline}

The pipeline processes 7 public datasets through four stages:

\begin{enumerate}
    \item \textbf{Download.} Fetch raw data from Hugging Face Hub, GitHub, and public APIs. Datasets range from 966 records (AlpacaEval) to 1.9M records (Open LLM Leaderboard).
    \item \textbf{Normalize.} Parse heterogeneous formats (CSV, JSON, Parquet) into a unified schema: \texttt{(model\_id, benchmark, metric, score, timestamp)}. Apply min-max normalization within each benchmark to produce $[0,1]$-bounded scores.
    \item \textbf{Map to Phi.} Route normalized records to the phi functions they inform, following a rule table that maps (dataset, metric) pairs to phi function IDs. For example, MMLU accuracy records map to $\phi_4$ (performance predictor) and $\phi_2$ (domain classifier), while pairwise comparison data maps to $\phi_8$ (cascade position).
    \item \textbf{Seed Weights.} Compute the initial weight vector $\wvec_0$ using Equation~\ref{eq:seed_weights}, combining literature priors with data quality multipliers.
\end{enumerate}

The pipeline is idempotent: re-running it on the same raw data produces identical outputs. Total processing time is under 5 minutes on a single machine.

\subsection{Normalized Database Statistics}

The normalized SQLite database contains 2,764,077 records across 15,526 unique model identifiers. Table~\ref{tab:phi_coverage} shows the data coverage per phi function.

\begin{table}[h]
\centering
\caption{Data coverage per phi function category.}
\label{tab:phi_coverage}
\small
\begin{tabular}{lcrrr}
\toprule
\textbf{Category} & \textbf{Count} & \textbf{Records} & \textbf{Models} & \textbf{$\sum w_i$} \\
\midrule
Data-grounded & 8 & 2,764,077 & 15,526 & 0.847 \\
Computed & 3 & 5,600 & 350 & 0.095 \\
Prior-based & 3 & 0 & 0 & 0.024 \\
Runtime/Derived & 2 & 0 & 0 & 0.024 \\
\midrule
\textbf{Total} & 16 & 2,764,077 & 15,526 & 1.000 \\
\bottomrule
\end{tabular}
\end{table}

The seed weight allocation reflects data quality: data-grounded phi functions receive 84.7\% of total weight, while prior-based and runtime functions share the remaining 15.3\%. This is a deliberate design choice, not a bug. The weight learner (Section~\ref{sec:convergence}) will adjust these weights as online feedback accumulates, potentially upweighting prior-based functions if they prove informative in practice.

\subsection{Seed Weight Computation}

The seed weights derive from two independent signals (Equation~\ref{eq:seed_weights}):

\paragraph{Literature Priors.} We surveyed three routing papers (FrugalGPT \citep{chen2023frugalgpt}, RouteLLM \citep{ong2024routellm}, and RouterBench \citep{hu2024routerbench}) and assigned importance scores on a 0.5--5.0 scale. Performance prediction ($\phi_4$, score 5.0) and historical success ($\phi_6$, score 4.0) consistently rank highest across all three papers. Cost efficiency ($\phi_{14}$, score 2.5) and instruction adherence ($\phi_{16}$, score 2.0) receive moderate ratings.

\paragraph{Data Quality Multipliers.} Each phi function receives a quality multiplier based on its data coverage:
\begin{itemize}
    \item High ($q=1.0$): $>$1,000 records and $>$50 models
    \item Medium ($q=0.5$): Moderate coverage
    \item Low ($q=0.25$): Limited coverage
    \item None ($q=0.1$): No direct data, prior-based only
\end{itemize}

The product $p_i \cdot q_i$, after simplex normalization, produces the seed weights reported in Table~\ref{tab:phi_functions}. The top-4 weights ($\phi_4 = 0.176$, $\phi_6 = 0.141$, $\phi_8 = 0.124$, $\phi_3 = \phi_{13} = 0.106$) together account for 65\% of total weight.

\subsection{Router Implementation}

The v0.1 router implements the full S(M,T) pipeline across 13 heterogeneous endpoints:

\begin{itemize}
    \item 7 LLM endpoints (Claude Opus 4, Claude Sonnet 4, GPT-4o, DeepSeek R1, DeepSeek V3, Qwen3 Coder, Gemini 2.5 Pro)
    \item 3 Agent endpoints (auto-search, mathematical verification, first-principles thinking)
    \item 1 Script endpoint (session capture)
    \item 2 MCP tool endpoints (Notion, Google Workspace)
\end{itemize}

This heterogeneous endpoint set distinguishes our implementation from prior routing work, which focuses exclusively on LLM-to-LLM routing. The gate layer handles type-specific constraints (agents have no context window limit, scripts have fixed capabilities, MCP tools have specific supported operations), while the phi functions provide unified scoring across all endpoint types.

\subsection{Smoke Test Results}

We verify correctness through four routing scenarios (Table~\ref{tab:smoke_tests}).

\begin{table}[h]
\centering
\caption{Smoke test results demonstrating correct routing behavior.}
\label{tab:smoke_tests}
\small
\begin{tabular}{p{4cm}lp{4cm}}
\toprule
\textbf{Query} & \textbf{Selected} & \textbf{Verification} \\
\midrule
``Write a Python sort function'' & DeepSeek V3 & Cheapest coding-capable model \\
``Solve $x^2 + 3x - 4 = 0$'' (thinking required) & DeepSeek R1 & 8/13 endpoints filtered for lacking thinking support \\
``Explain quantum entanglement'' & Claude Sonnet 4 & Top general-knowledge model \\
List endpoints & 13 registered & 7 LLM + 3 agent + 1 script + 2 MCP \\
\bottomrule
\end{tabular}
\end{table}

These tests verify gate filtering (thinking requirement eliminates 8 endpoints), phi function evaluation (coding task favors code-optimized models), and cost sensitivity (DeepSeek V3 selected over more expensive alternatives for a straightforward coding task).


\section{Learning Phi Functions: The BilinearPhiEngine}
\label{sec:learned_phi}

Section~\ref{sec:measurement} established that hand-designed phi functions cannot be grounded in public benchmark data due to structural problems: metric incomparability, domain transfer breakdown, and dimensional collapse. This section presents an alternative: learning the phi functions end-to-end from pairwise routing data, bypassing the measurement gap entirely.

\subsection{Motivation}

The measurement gap arises because hand-designed phi functions impose semantic structure (``domain classifier,'' ``reasoning depth'') that requires mapping heterogeneous benchmark metrics into a common space. This mapping is where all 9 designs failed. But routing does not require interpretable phi functions. It requires phi functions that \emph{predict routing outcomes}. If we replace the 16 semantic phis with 16 learned bilinear functions trained directly on routing-relevant data, the normalization and domain transfer problems disappear: the model learns whatever internal representation best predicts scores, without requiring us to specify what each dimension means.

Learned representations for routing are not new. RouteLLM \citep{ong2024routellm} uses matrix factorization over preference data. RouterDC \citep{chen2024routerdc} trains dual contrastive embeddings. EmbedLLM \citep{zhuang2024embedllm} learns compact model representations for performance prediction. GraphRouter \citep{feng2024graphrouter} constructs heterogeneous task-model graphs. Our contribution here is not the bilinear form itself, but its role as a drop-in replacement for the hand-designed phi functions within the S(M,T) architecture: preserving the framework's multiplicative gating and convergence properties while bypassing the measurement gap.

\subsection{Architecture}

The BilinearPhiEngine instantiates the weighted compatibility term of S(M,T) as a set of learned bilinear forms:

\begin{equation}
\label{eq:bilinear_phi}
\hat{\phi}_i(M, T) = \sigma\!\left(\mathbf{e}_T^\top \mathbf{W}_i \, \mathbf{e}_M + b_i\right), \quad i = 1, \ldots, K
\end{equation}

where $\mathbf{e}_T \in \R^{768}$ is a frozen prompt embedding from a pre-trained sentence encoder (all-mpnet-base-v2), $\mathbf{e}_M \in \R^{256}$ is a learned model embedding, $\mathbf{W}_i \in \R^{768 \times 256}$ is a learned bilinear interaction matrix, $b_i \in \R$ is a bias term, and $\sigma$ is the sigmoid function. The overall score combines these via softmax-weighted sum:

\begin{equation}
\label{eq:bilinear_score}
\hat{S}(M, T) = \sum_{i=1}^{K} \alpha_i \cdot \hat{\phi}_i(M, T), \quad \alpha_i = \frac{e^{a_i}}{\sum_j e^{a_j}}
\end{equation}

where $\mathbf{a} \in \R^K$ are learned logits producing weights on the probability simplex. With $K=16$ bilinear heads, the architecture has approximately 3.15 million trainable parameters:

\begin{itemize}
    \item 16 bilinear matrices $\mathbf{W}_i$: $16 \times 768 \times 256 = 3{,}145{,}728$ parameters
    \item 16 biases $b_i$: 16 parameters
    \item 48 model embeddings $\mathbf{e}_M$: $48 \times 256 = 12{,}288$ parameters
    \item 16 head weight logits $a_i$: 16 parameters
\end{itemize}

The prompt encoder is frozen during training. Only model embeddings, bilinear matrices, head weights, and biases are learned. This separation is deliberate: the prompt encoder provides a fixed, high-quality text representation; the bilinear matrices learn how different models interact with different prompt features.

\paragraph{Diversity Regularization.} To prevent the 16 heads from collapsing to identical functions, we add a cosine-similarity penalty between head weight vectors:

\begin{equation}
\label{eq:diversity}
\mathcal{L}_{\text{div}} = \frac{1}{K(K-1)} \sum_{i \neq j} \left|\cos(\text{vec}(\mathbf{W}_i), \text{vec}(\mathbf{W}_j))\right|
\end{equation}

where $\text{vec}(\mathbf{W}_i)$ flattens the bilinear matrix into a vector. The total loss is $\mathcal{L} = \text{MSE}(\hat{S}, S^*) + \lambda_{\text{div}} \cdot \mathcal{L}_{\text{div}}$ with $\lambda_{\text{div}} = 0.01$.

\subsection{Training Data}

We unify five public routing datasets into a single training corpus of 740,803 (prompt, model, score) triples:

\begin{table}[h]
\centering
\caption{Training data sources with canonical model resolution.}
\label{tab:training_data}
\small
\begin{tabular}{lrrl}
\toprule
\textbf{Source} & \textbf{Samples} & \textbf{Models} & \textbf{Score Type} \\
\midrule
RouterBench & 401,467 & 11 & Accuracy (0/1) \\
RouteLLM & 218,202 & 2 & Win rate (0/1) \\
Arena 55K & 111,712 & 45 & Elo-derived (continuous) \\
MT-Bench & 6,710 & 5 & Rating (1--10, normalized) \\
RewardBench & 2,712 & 31 & Accuracy (0/1) \\
\midrule
\textbf{Total} & \textbf{740,803} & \textbf{48} & \\
\bottomrule
\end{tabular}
\end{table}

A canonical model registry maps 180 model name aliases (e.g., ``gpt-4-1106-preview,'' ``gpt4-turbo,'' ``openai/gpt-4-1106'') to 48 canonical model identifiers, each receiving a learned 256-dimensional embedding. Scores are normalized to $[0, 1]$ per source. Data is split 80/10/10 into train/validation/test sets with a fixed random seed.

\paragraph{Prompt Encoding.} The 162,422 unique prompts are pre-encoded using the raw transformer model (AutoTokenizer + AutoModel from all-mpnet-base-v2) with mean pooling and L2 normalization, producing 768-dimensional embeddings. This one-time encoding step takes approximately 45 minutes on Apple Silicon (MPS); the resulting cache enables the training loop to run at full GPU speed without I/O bottleneck.

\subsection{Training Procedure}

Training uses AdamW with differential learning rates: $10^{-4}$ for bilinear matrices $\mathbf{W}_i$, $10^{-3}$ for model embeddings $\mathbf{e}_M$ and head weights $\mathbf{a}$. Batch size is 256. Early stopping with patience 10 monitors validation MSE loss.

The full model (all 5 sources) trains for 31 epochs before early stopping, reaching best validation loss of 0.1217. Training takes approximately 2 minutes per epoch on Apple M-series hardware (MPS backend), for a total wall time under 70 minutes.

\subsection{Benchmark Evaluation}

We evaluate the trained BilinearPhiEngine on two established routing benchmarks. A third benchmark (RouterEval) proved incompatible because it evaluates routing among 1,000+ open-source fine-tunes from the Open LLM Leaderboard, a model pool completely disjoint from our 48 API models.

\subsubsection{RouterBench}

RouterBench \citep{hu2024routerbench} provides 73,008 test instances across 11 model pairs, split into 0-shot and 5-shot settings. For each instance, the router must select the model most likely to answer correctly.

\begin{table}[h]
\centering
\caption{RouterBench evaluation results.}
\label{tab:routerbench_results}
\small
\begin{tabular}{lccc}
\toprule
\textbf{Method} & \textbf{0-shot Acc.} & \textbf{5-shot Acc.} & \textbf{Overall} \\
\midrule
Random baseline & 55.50 & 64.34 & -- \\
Always-strong & 84.35 & 86.66 & 85.51 \\
\textbf{BilinearPhi (ours)} & \textbf{80.49} & \textbf{86.76} & \textbf{83.63} \\
Oracle (upper bound) & 96.13 & 96.76 & 96.45 \\
\bottomrule
\end{tabular}
\end{table}

The model achieves 83.63\% overall accuracy. In the 5-shot setting, it slightly exceeds the always-strong baseline (86.76\% vs. 86.66\%), indicating it learns when cheaper models suffice. The router selects from 11 distinct models, not just the strongest, demonstrating genuine multi-model routing rather than defaulting to a single endpoint.

\subsubsection{RouteLLM}

RouteLLM \citep{ong2024routellm} evaluates binary routing (strong vs. weak model) across MMLU, GSM8K, and MT-Bench subsets. Quality is measured as the average win rate of selected models, with AUC computed across routing thresholds.

\begin{table}[h]
\centering
\caption{RouteLLM evaluation results. Quality ratio = quality / always-strong quality.}
\label{tab:routellm_results}
\small
\begin{tabular}{lcc}
\toprule
\textbf{Threshold} & \textbf{Quality} & \textbf{Quality Ratio} \\
\midrule
Always weak & 0.7187 & 84.09\% \\
20\% strong & 0.7609 & 89.02\% \\
50\% strong & 0.8065 & 94.35\% \\
80\% strong & 0.8399 & 98.27\% \\
Always strong & 0.8547 & 100.00\% \\
\midrule
\textbf{AUC} & \multicolumn{2}{c}{\textbf{0.8006}} \\
Routing accuracy & \multicolumn{2}{c}{94.55\%} \\
\bottomrule
\end{tabular}
\end{table}

At the 50\% strong-model threshold (routing half of queries to the cheaper model), the system retains 94.35\% of always-strong quality. The AUC of 0.8006 indicates effective quality-cost tradeoff across the full threshold range.

\subsubsection{Internal Evaluation}

On the held-out test split (10\% of training data, 74,080 samples), the model achieves MSE 0.1223, MAE 0.2689, Pearson correlation 0.5305, and binary routing accuracy 74.39\%.

\subsection{Leave-One-Out Generalization Analysis}

To test whether the BilinearPhiEngine learns transferable routing features rather than source-specific patterns, we conduct a leave-one-out ablation study. Five models are trained, each excluding one of the five data sources entirely. Each ablated model is then evaluated on its held-out source (zero-shot generalization) and on all five sources (cross-source transfer).

\subsubsection{Training Dynamics}

\begin{table}[h]
\centering
\caption{Leave-one-out ablation: training summary. Correlation is Pearson $r$ on the held-out source.}
\label{tab:loo_training}
\small
\begin{tabular}{lrrrrr}
\toprule
\textbf{Excluded} & \textbf{Train $N$} & \textbf{Epochs} & \textbf{Val Loss} & \textbf{Heldout $r$} & \textbf{Heldout MSE} \\
\midrule
RouterBench (401K) & 271,470 & 21$^\dagger$ & 0.0752 & 0.6505 & 0.0790 \\
RouteLLM (218K) & 418,081 & 17$^\dagger$ & 0.1623 & 0.3192 & 0.1742 \\
Arena 55K (112K) & 503,273 & 30 & 0.1131 & 0.4596 & 0.1794 \\
MT-Bench (7K) & 587,275 & 30 & 0.1211 & 0.4690 & 0.1705 \\
RewardBench (3K) & 590,473 & 30 & 0.1205 & 0.3849 & 0.1539 \\
\midrule
Full model & 592,641 & 31$^\dagger$ & 0.1217 & -- & 0.1223 \\
\bottomrule
\end{tabular}
\begin{flushleft}
\small $^\dagger$ Early stopped (patience 8).
\end{flushleft}
\end{table}

Two patterns stand out. First, removing RouterBench (54\% of data) causes early stopping at epoch 21, while removing smaller sources allows training to run all 30 epochs. This confirms RouterBench's role as the diversity backbone of the dataset. Second, removing MT-Bench (0.9\% of data) or RewardBench (0.4\%) produces validation loss nearly identical to the full model (0.1211 and 0.1205 vs. 0.1217), indicating these sources add negligible marginal information.

\subsubsection{Cross-Source Generalization Matrix}

Table~\ref{tab:cross_source} presents the full 5$\times$5 cross-source correlation matrix. Each row is a model trained \emph{without} that source; each column is the evaluation target. Diagonal entries (bold) represent zero-shot generalization to an unseen source.

\begin{table}[h]
\centering
\caption{Cross-source Pearson correlation matrix. Rows: excluded source during training. Columns: evaluation source. Diagonal (bold) = zero-shot generalization.}
\label{tab:cross_source}
\small
\begin{tabular}{lccccc}
\toprule
\textbf{Trained w/o} & \textbf{RBench} & \textbf{RLLM} & \textbf{Arena} & \textbf{MTB} & \textbf{RwdB} \\
\midrule
RouterBench & \textbf{0.651} & 0.073 & 0.197 & 0.174 & 0.169 \\
RouteLLM & 0.425 & \textbf{0.319} & 0.457 & 0.468 & 0.385 \\
Arena 55K & 0.583 & 0.321 & \textbf{0.460} & 0.470 & 0.386 \\
MT-Bench & 0.682 & 0.321 & 0.459 & \textbf{0.469} & 0.385 \\
RewardBench & 0.679 & 0.320 & 0.459 & 0.468 & \textbf{0.385} \\
\bottomrule
\end{tabular}
\end{table}

\paragraph{Finding 1: RouterBench is the generalization anchor.}
When RouterBench is included in training (rows 2--5), the RouterBench column shows correlation 0.43--0.68. When RouterBench is excluded (row 1), cross-source generalization collapses: correlation on all non-RouterBench sources falls to 0.07--0.20. RouterBench's 11-model, diverse-prompt format provides the signal diversity that enables generalization.

\paragraph{Finding 2: RouteLLM is the hardest transfer target.}
The RouteLLM column is capped at $r \approx 0.32$ regardless of which other source is removed. RouteLLM's binary 2-model setup (GPT-4 vs. Mixtral on academic benchmarks) represents a qualitatively different routing problem that other benchmark formats cannot fully predict.

\paragraph{Finding 3: Small sources are redundant.}
Rows 4--5 (excluding MT-Bench or RewardBench) produce nearly identical cross-source correlations to each other and to the full model. Training on just RouterBench + RouteLLM + Arena 55K (731K samples) achieves the same generalization quality as the full 741K dataset.

\paragraph{Finding 4: Unseen sources are predictable.}
The diagonal entries (zero-shot generalization) range from 0.319 to 0.651, confirming that the bilinear architecture learns transferable routing features. RouterBench is the most predictable from other sources ($r = 0.651$), while RouteLLM is the least ($r = 0.319$).

\subsection{Connection to the Measurement Gap}

The BilinearPhiEngine achieves what the hand-designed phi functions could not: reliable routing scores grounded in public benchmark data. The resolution is architectural. Hand-designed phis require explicit metric normalization, domain classification, and semantic mapping, each of which introduces fragile assumptions (Section~\ref{sec:measurement}). The bilinear architecture sidesteps all three by learning the mapping end-to-end. The prompt encoder provides a fixed text representation; the bilinear matrices learn whatever model-prompt interaction structure best predicts routing outcomes.

The tradeoff is interpretability. The 16 learned bilinear heads have no semantic labels. We cannot say ``head 3 captures reasoning depth'' or ``head 7 measures instruction adherence.'' The hand-designed phi functions in Table~\ref{tab:phi_functions} remain the interpretable interface for the production router (Section~\ref{sec:empirical}), initialized with seed weights and refined through online learning. The bilinear engine provides the empirical validation that the weighted-phi architecture can achieve competitive routing accuracy when the learning problem is well-posed.

\subsection{Scaling Ablation: BilinearPhiEngine V3}
\label{sec:v3_ablation}

The BilinearPhiEngine V1 (Section~\ref{sec:learned_phi}) was trained on a single data corpus with a compact architecture (3.15M parameters, $K=16$ heads, $D_{\text{model}}=256$). A natural question is whether scaling the architecture with more data, more models, and additional inductive biases improves routing quality. We designed BilinearPhiEngine V3 to test this hypothesis. The answer is no: V3 underperforms V1 on every benchmark.

\subsubsection{V3 Architecture and Training}

V3 scales three dimensions of V1 simultaneously:

\begin{enumerate}
    \item \textbf{Parameter expansion.} $K=24$ bilinear heads (vs.\ 16), $D_{\text{model}}=768$ (vs.\ 256), producing bilinear matrices $\mathbf{W}_i \in \R^{768 \times 768}$. Total parameters: 14.6M (4.63$\times$ V1).
    \item \textbf{Data expansion.} Training on a V2 database (4.0 GB, 3.86M rows from 10 additional preference datasets including UltraFeedback, Nectar, WildBench, HelpSteer2, and HH-RLHF). After filtering unrecognized model IDs (53 of 101 models), 1.98M usable rows with 356,358 unique prompts remain (2.2$\times$ the unique prompts in V1).
    \item \textbf{Architectural additions.} Three new components added via progressive training stages: task-type conditioning embeddings (Stage~3), multi-head cross-attention between prompt and model representations (Stage~4), and metadata feature integration including source and score normalization information (Stage~5).
\end{enumerate}

Training uses a 5-stage progressive schedule on an A100 80GB GPU: (1) frozen prompt encoder with bilinear matrices only, (2) full end-to-end fine-tuning including unfreezing the sentence encoder, (3) task conditioning, (4) cross-attention, (5) metadata features. Each stage initializes from the previous stage's best checkpoint.

\subsubsection{Results}

We evaluate the two most informative checkpoints: Stage~2 (the first fully trained model, before architectural additions) and Stage~5 (the final model with all additions).

\paragraph{External Benchmarks.}

\begin{table}[h]
\centering
\caption{V1 vs.\ V3 on external routing benchmarks. V1 outperforms all V3 variants. $\Delta$ is percentage points relative to V1.}
\label{tab:v3_external}
\small
\begin{tabular}{lcccccc}
\toprule
& \multicolumn{3}{c}{\textbf{RouterBench}} & \multicolumn{3}{c}{\textbf{RouteLLM}} \\
\cmidrule(lr){2-4} \cmidrule(lr){5-7}
\textbf{Model} & \textbf{0-shot} & \textbf{5-shot} & \textbf{Overall} & \textbf{AUC} & \textbf{Q@50\%} & \textbf{Rt.\ Acc.} \\
\midrule
\textbf{V1} (3.15M) & \textbf{80.49} & \textbf{86.76} & \textbf{83.63} & \textbf{0.8006} & \textbf{0.9435} & \textbf{94.55} \\
V3-S2 (14.6M) & 79.18 & 86.70 & 82.94 & 0.7972 & 0.8020 & 93.10 \\
V3-S5 (14.6M) & 77.44 & 85.51 & 81.47 & 0.7973 & 0.8022 & 92.79 \\
\midrule
$\Delta$ (best V3) & $-1.31$ & $-0.06$ & $-0.69$ & $-0.0033$ & $-0.1413$ & $-1.45$ \\
$\Delta$ (worst V3) & $-3.05$ & $-1.25$ & $-2.16$ & $-0.0034$ & $-0.1415$ & $-1.76$ \\
\bottomrule
\end{tabular}
\end{table}

On RouterBench, V3 Stage~2 loses 0.69 percentage points overall, while Stage~5 loses 2.16 points. On RouteLLM, both V3 stages lose approximately 0.003 AUC. The quality-at-50\%-strong metric shows a larger gap: V3 retains only 80.2\% of always-strong quality at 50\% cost (vs.\ V1's 94.4\%), suggesting degraded confidence calibration.

\paragraph{Task-Stratified Analysis.}

\begin{table}[h]
\centering
\caption{RouterBench accuracy by task type. V3-S2 is the best V3 variant. Bold indicates the best result per task.}
\label{tab:v3_task}
\small
\begin{tabular}{lrrr}
\toprule
\textbf{Task Type} & \textbf{V3-S2 (\%)} & \textbf{V3-S5 (\%)} & \textbf{$N$} \\
\midrule
Math & 94.39 & \textbf{94.37} & 15{,}547 \\
Reasoning & 81.12 & 79.03 & 44{,}922 \\
Safety & 79.79 & 79.04 & 1{,}064 \\
Code & 77.80 & \textbf{78.00} & 500 \\
Factual QA & 77.95 & 76.65 & 6{,}616 \\
Instr.\ following & 75.55 & 73.36 & 687 \\
Conversation & 72.88 & 70.77 & 520 \\
Writing & 68.05 & 67.67 & 3{,}152 \\
\bottomrule
\end{tabular}
\end{table}

V3 is competitive on math tasks (94.4\%), which have unambiguous correctness labels. Performance degrades most on tasks requiring subjective quality judgments: writing (68\%), conversation (71--73\%), and instruction following (74--76\%). This pattern is consistent with the V2 training data containing noisier labels from preference-based datasets.

\paragraph{Internal Evaluation.}

On the held-out V2 test split (198,304 samples), V3 shows near-zero or negative correlation with ground truth:

\begin{table}[h]
\centering
\caption{Internal evaluation on V2 test data. V1 baseline computed on V1 test split.}
\label{tab:v3_internal}
\small
\begin{tabular}{lcccc}
\toprule
\textbf{Model} & \textbf{Pearson $r$} & \textbf{MSE} & \textbf{MAE} & \textbf{Binary Acc.\ (\%)} \\
\midrule
V1 & 0.5305 & 0.1223 & 0.2689 & 74.39 \\
V3-S2 & $-0.0483$ & 0.1365 & 0.3188 & 57.99 \\
V3-S5 & $-0.0335$ & 0.1357 & 0.3162 & 58.09 \\
\bottomrule
\end{tabular}
\end{table}

The internal correlation collapse ($r = 0.53 \to r \approx -0.04$) is more dramatic than the external benchmark regressions. This divergence has a structural explanation: the V2 database labels incorporate preference-based scores from datasets with different quality semantics than V1's benchmark-accuracy scores. V3 learned the V2 label distribution, but that distribution is not well-aligned with the actual model quality ranking.

\subsubsection{Why Scaling Failed: Analysis}

Four factors explain V3's underperformance:

\paragraph{1. Label distribution shift.} The V2 database merged 10 new preference datasets with the original 5 benchmark sources. Of 101 model IDs in V2, 53 were unrecognized by the canonical registry, causing 1.87M rows (49\% of total) to be skipped. The surviving data has different score semantics: preference-based ``win rate'' labels from UltraFeedback and Nectar encode relative quality between model pairs, not absolute task accuracy. Training on these labels teaches V3 to predict pairwise preferences rather than absolute routing scores.

\paragraph{2. Sparse model coverage.} V3's embedding table covers 93 models (vs.\ V1's 48), but the new models have far fewer training examples. The V2 database has 356K unique prompts spread across 93 models, versus V1's 162K prompts across 48 models. Per-model embedding quality degrades with sparse coverage, and the embedding space becomes undertrained for tail models.

\paragraph{3. Progressive training instability.} Each stage inherits and potentially amplifies errors from its predecessor. Stage~5 (the final model) performs worse than Stage~2 (the first fully trained model) on RouterBench ($81.47\%$ vs.\ $82.94\%$), confirming that the architectural additions (task conditioning, cross-attention, metadata) degraded rather than improved routing quality. This is characteristic of overfitting to training artifacts rather than learning generalizable routing features.

\paragraph{4. Architectural overparameterization.} With 14.6M parameters trained on 1.98M usable samples, V3 has a parameter-to-sample ratio of 7.4:1. V1's ratio is 4.3:1 with cleaner labels. The bilinear matrices $\mathbf{W}_i \in \R^{768 \times 768}$ have 590K parameters each (vs.\ V1's 197K), providing excess capacity that absorbed noise in the V2 labels rather than learning better routing representations.

\subsubsection{Implications}

The V3 ablation establishes that the original V1 architecture is not capacity-limited. The bilinear routing problem, at least at the current scale of 48--93 models, does not benefit from larger embedding dimensions, more bilinear heads, or additional architectural components. V1's 3.15M parameters with clean labels outperform V3's 14.6M parameters with noisy labels on every metric. This result has two practical implications.

First, \textbf{data quality dominates model capacity} for learned routing. The path to better routing accuracy runs through cleaner, more consistent training labels, not larger models. This aligns with the measurement gap finding (Section~\ref{sec:measurement}): the bottleneck is grounding quality signals, not the capacity of the model consuming those signals.

Second, \textbf{progressive architectural additions provide diminishing returns} when the base training signal is noisy. Task conditioning, cross-attention, and metadata features all assume the training labels contain structured information that can be decomposed along these axes. When labels are inconsistent, these additions fit noise rather than signal.

The V3 experiment thus serves as a controlled counter-test to the V1 results. V1's performance is not an accident of small-scale evaluation. A systematic attempt to improve it through data scaling, parameter scaling, and architectural enrichment produced consistently worse results. The compact bilinear architecture, trained on well-curated data, remains the strongest configuration.

\section{Related Work}
\label{sec:related}

\paragraph{LLM Routing and Cascading.}
FrugalGPT \citep{chen2023frugalgpt} pioneered cost-aware LLM selection through cascading strategies, achieving 50--90\% cost reduction. Hybrid LLM \citep{ding2024hybridllm} routes queries between small and large models based on predicted difficulty, reducing large-model calls by up to 40\%. RouteLLM \citep{ong2024routellm} trained pairwise classifiers for binary routing using preference data from LMSYS Chatbot Arena. AutoMix \citep{madaan2023automix} used self-verification to decide when to escalate from small to large models. \citet{dekoninck2024cascade} unified routing and cascading into a single framework, showing that iterative model selection with optional skipping outperforms pure routing or pure cascading by up to 14\%. BEST-Route \citep{ding2025bestroute} extended this by jointly selecting both model and number of responses based on query difficulty. RouterBench \citep{hu2024routerbench} and RouterEval \citep{huang2024routereval} provide standardized evaluation frameworks. Most of these systems focus on LLM-to-LLM routing with additive or sequential scoring.

\paragraph{Learned Representations for Routing.}
A growing body of work learns query and model representations for routing decisions. RouteLLM's strongest baseline uses matrix factorization over preference data. RouterDC \citep{chen2024routerdc} trains a router via dual contrastive learning (sample-LLM and sample-sample losses) to match queries with the best-performing LLM from a candidate pool. EmbedLLM \citep{zhuang2024embedllm} learns compact vector embeddings of LLMs using an encoder-decoder approach, enabling performance forecasting across 112+ models. GraphRouter \citep{feng2024graphrouter} constructs heterogeneous task-model graphs and uses edge prediction for routing. ICL-Router \citep{wang2025iclrouter} uses in-context vectors to represent model capabilities, enabling integration of new models without retraining. Prompt-to-Leaderboard \citep{frick2025p2l} trains an LLM to output Bradley-Terry coefficients from prompts, predicting human preferences for prompt-specific model ranking. ZeroRouter \citep{yan2026zerorouter} creates a universal, model-agnostic latent space for query difficulty, enabling zero-shot onboarding of new models using only anchor queries. Our BilinearPhiEngine (Section~\ref{sec:learned_phi}) shares the learned-embedding approach with these systems. The specific contribution is integrating bilinear forms as drop-in replacements for the hand-designed phi functions within the S(M,T) scoring architecture, preserving the framework's multiplicative gating and convergence properties while bypassing the measurement gap.

\paragraph{Multi-Agent and Heterogeneous Routing.}
Several recent systems extend routing beyond LLM-to-LLM selection. MoMA \citep{moma2025routing} integrates LLM-as-a-judge with Mixture-of-Experts for unified routing across both models and agents. DAAO \citep{daao2025} uses a VAE for difficulty estimation and dynamically generates query-specific multi-agent workflows. xRouter \citep{qian2025xrouter} frames routing as a tool-calling problem and trains end-to-end with reinforcement learning. Router-R1 \citep{zhang2025routerr1} teaches an LLM to interleave reasoning and routing actions for multi-round aggregation. S(M,T) also routes across heterogeneous endpoints (LLMs, agents, scripts, tool-use systems), but through a different mechanism: a single scoring equation with endpoint-type-specific gates rather than an LLM-based orchestrator.

\paragraph{Mixture of Experts.}
Sparse Mixture-of-Experts \citep{shazeer2017outrageously} routes tokens to specialized sub-networks within a single model. Switch Transformers \citep{fedus2022switch} simplified gating to top-1 expert selection. These approaches operate at the token level within a model, while S(M,T) operates at the query level across models. The Product-of-Experts framework \citep{hinton2002poe}, which we use for the gate layer, was originally proposed for combining generative models.

\paragraph{Online Learning for Model Selection.}
The contextual bandit formulation for model selection has been explored by \citet{agrawal2013thompson} in the general setting. TI-UCB \citep{tiucb2024} adapts upper confidence bound methods for online model selection with time-increasing reward signals. BaRP \citep{barp2025} frames routing as a multi-objective contextual bandit conditioned on trade-off vectors, enabling performance-cost preference adjustment at inference without retraining. OmniRouter \citep{mei2025omnirouter} models routing as constrained optimization via Lagrangian dual decomposition under budget constraints. Robbins-Monro stochastic approximation \citep{robbins1951stochastic} and convergence under Hajek conditions \citep{hajek1988cooling} are well-established. Soft Actor-Critic \citep{haarnoja2018sac} introduced entropy-regularized learning for continuous control. We combine Robbins-Monro weight updates with SAC-style entropy decay for routing, with Duchi simplex projection \citep{duchi2008simplex} to maintain valid weight distributions.

\paragraph{Benchmark Reliability.}
Concerns about benchmark contamination \citep{deng2023leaderboards}, evaluation methodology \citep{liang2022helm}, and the gap between benchmark performance and real-world utility have been raised repeatedly. Our measurement gap analysis (Section~\ref{sec:measurement}) provides a systematic documentation of how these concerns specifically impact routing, demonstrating 0/14 successful measurement designs across 2.76M records.

\paragraph{What S(M,T) Adds.}
Given the breadth of recent routing work, we clarify the specific contributions of this paper. (1) \emph{A unified scoring equation with formal guarantees}: the three-part multiplicative form (gate products $\times$ weighted compatibility $\times$ Boltzmann cost) with provable convergence ($O(\sqrt{KN\log N})$ regret), parameter identifiability, and adversarial robustness. No prior system provides all three guarantees for the complete routing pipeline. (2) \emph{The measurement gap as an empirical finding}: showing that hand-designed, interpretable scoring functions cannot be grounded in public benchmark data due to metric incomparability, domain transfer breakdown, and dimensional collapse. This finding is independent of any particular routing architecture. (3) \emph{Heterogeneous endpoint support via type-specific gates}: while MoMA and DAAO also route beyond LLMs, S(M,T) handles heterogeneous endpoints through the gate layer within a single equation rather than through an LLM orchestrator, preserving scoring interpretability and convergence properties. (4) \emph{A learned resolution that preserves framework structure}: the BilinearPhiEngine replaces hand-designed phi functions with learned bilinear forms within the same S(M,T) architecture. The bilinear form itself builds on prior work in learned routing representations (RouterDC, EmbedLLM, GraphRouter). The contribution is demonstrating that this approach bypasses the measurement gap while retaining the framework's multiplicative gating and weight-learning machinery.

\section{Discussion}
\label{sec:discussion}

\subsection{The Interpretability-Performance Tradeoff}

The most striking finding of this work is a tradeoff, not an impossibility. All 14 theoretical results (convergence, regret bounds, identifiability, orthogonality, adversarial robustness) verified at 85--95\% confidence. Zero of 14 hand-designed measurement designs survived stress-testing. But the BilinearPhiEngine, using 16 learned bilinear forms on the same benchmark data, achieves 83.63\% routing accuracy on RouterBench and AUC 0.8006 on RouteLLM.

The gap separates \emph{interpretable} from \emph{learned} routing. Hand-designed phi functions require explicit metric normalization, domain classification, and semantic mapping. Each of these steps introduces fragile assumptions that benchmark data cannot support. Learned bilinear forms sidestep all three by training end-to-end. The cost is that the 16 learned heads carry no semantic labels. We cannot say ``head 3 captures reasoning depth'' or ``head 7 measures instruction adherence.''

The V3 scaling ablation (Section~\ref{sec:v3_ablation}) strengthens this finding. A 4.63$\times$ parameter increase with 2.2$\times$ more unique prompts and three architectural additions (task conditioning, cross-attention, metadata) uniformly degraded performance. The bottleneck is not model capacity but label quality: preference-based scores from heterogeneous datasets encode different semantics than benchmark accuracy, and naive data scaling amplifies this inconsistency rather than averaging it out.

This tradeoff has practical implications. For production routing where operators need to audit, debug, and explain routing decisions, the hand-designed phi functions with literature-informed seed weights and online learning remain the appropriate interface. For empirical validation that the weighted-phi architecture \emph{can} achieve competitive accuracy, the bilinear engine provides the evidence. The two serve complementary roles.

\subsection{Limitations}

\paragraph{Seed Weight Subjectivity.} The literature prior scores $p_i$ (Section~\ref{sec:empirical}) were assigned by a single researcher surveying three papers. Different reviewers might assign different scores, shifting the seed weight vector. The Robbins-Monro learning procedure corrects for initialization bias over time (Proposition~\ref{prop:prior}), but early routing decisions are sensitive to the prior.

\paragraph{Small Endpoint Set.} The v0.1 router covers 13 endpoints. Production routing systems may handle hundreds. The $O(\sqrt{KN\log N})$ regret bound scales with $K$ (phi dimensions, not endpoint count), but the gate evaluation cost scales linearly with the endpoint set size.

\paragraph{No A/B Testing.} We have not compared S(M,T) routing decisions against human expert routing or competing algorithms on live traffic. The smoke tests (Table~\ref{tab:smoke_tests}) verify correctness of the pipeline, not optimality of the routing policy.

\paragraph{Single-Turn Focus.} The current framework scores each task independently. Multi-turn conversations, where context accumulates and the optimal endpoint may change mid-conversation, are not addressed.

\paragraph{Static Phi Definitions.} The 16 phi functions are defined at design time. New capability dimensions (e.g., multimodal reasoning, tool-use orchestration) would require adding new phi functions and re-deriving seed weights.

\subsection{Broader Impact}

LLM routing affects cost and quality for anyone using language models. A good router reduces cost by directing simple tasks to cheap models while preserving quality on hard tasks. A bad router wastes money on overqualified models or produces poor results by underestimating task difficulty. The measurement gap finding has a specific implication: routing systems that claim benchmark-grounded accuracy should be scrutinized carefully. The normalization decisions, domain transfer assumptions, and contamination risks documented in Section~\ref{sec:measurement} apply to all benchmark-based routing approaches, not just ours.

\subsection{Future Directions}

\paragraph{Routing-Specific Benchmarks.} The field needs benchmarks designed for routing evaluation, not repurposed from model comparison. These would measure task-conditional, model-specific performance under deployment conditions, with controlled contamination and standardized metrics.

\paragraph{Multi-Turn Routing.} Extending S(M,T) to multi-turn conversations requires a state variable tracking accumulated context, cost, and user satisfaction across turns. The scoring equation becomes $S(M, T, s_t)$ where $s_t$ is the conversation state at turn $t$.

\paragraph{Federated Weight Learning.} Multiple router deployments could share learned weights (after differential privacy protection) to accelerate convergence. The Robbins-Monro framework supports aggregated gradient estimates from distributed sources.

\paragraph{Bridging Interpretability and Performance.} The BilinearPhiEngine (Section~\ref{sec:learned_phi}) demonstrates that learned phi functions achieve competitive routing accuracy. The open question is whether the learned bilinear heads can be \emph{interpreted} post-hoc, perhaps through probing classifiers or attention-style analysis, to recover semantic labels. If so, the interpretability-performance tradeoff could be narrowed. The diversity regularization already encourages the heads to capture different aspects of model-task interaction; understanding \emph{what} each head captures would connect the learned and hand-designed approaches.

\paragraph{Data Curation Over Model Scaling.} The V3 ablation (Section~\ref{sec:v3_ablation}) suggests that future routing improvements should prioritize training data curation over architectural scaling. Specifically: consistent scoring semantics across sources, rigorous model identity resolution, and filtering for label quality before increasing training set size. This echoes findings in the inverse scaling literature \citep{mckenzie2023inverse}, where more capacity on noisy signals degrades performance.

\section{Conclusion}
\label{sec:conclusion}

We presented S(M,T), a unified scoring framework for routing queries to optimal endpoints. The framework combines multiplicative gating, weighted compatibility functions, and Boltzmann cost penalties into a single equation with provable convergence ($O(\sqrt{KN\log N})$ regret), parameter identifiability, and adversarial robustness. It generalizes beyond LLM-to-LLM routing to cover agents, scripts, and tool-use systems through a common scoring interface.

Three empirical findings define the paper's contribution. First, the measurement gap: all 14 attempts to ground hand-designed phi functions in 2.76 million normalized benchmark records failed due to structural problems (metric incomparability, domain transfer breakdown, dimensional collapse). Second, the learned resolution: replacing the 16 hand-designed phis with 16 learned bilinear forms trained on 740,803 pairwise routing samples achieves 83.63\% accuracy on RouterBench and AUC 0.8006 on RouteLLM, retaining 94.35\% of always-strong quality at 50\% cost reduction. Third, a scaling ablation: a 4.63$\times$ larger model (14.6M parameters) trained on expanded data with progressive architectural additions underperforms the compact V1 on every benchmark, establishing that data quality dominates model capacity for learned routing.

The measurement gap thus separates interpretable from learned routing, not theory from practice. Hand-designed phi functions that an operator can audit and explain cannot be grounded in public benchmark data. Learned bilinear forms that carry no semantic labels can. The tradeoff between interpretability and performance is the central tension this paper identifies, and narrowing that tradeoff is the most promising direction for future work.

The v0.1 router is operational, routing across 13 heterogeneous endpoints with literature-informed seed weights for the interpretable interface and bilinear heads for empirical validation.

\paragraph{AI Disclosure.} AI tools developed by Anthropic were used to assist with writing, code development, and experimental analysis in the preparation of this work. All research direction, experimental design, and intellectual contributions are the author's own.


% === Bibliography ===
\bibliography{bibliography}

\end{document}
